% Options for packages loaded elsewhere
\PassOptionsToPackage{unicode}{hyperref}
\PassOptionsToPackage{hyphens}{url}
\documentclass[
]{article}
\usepackage{xcolor}
\usepackage[margin=1in]{geometry}
\usepackage{amsmath,amssymb}
\setcounter{secnumdepth}{-\maxdimen} % remove section numbering
\usepackage{iftex}
\ifPDFTeX
  \usepackage[T1]{fontenc}
  \usepackage[utf8]{inputenc}
  \usepackage{textcomp} % provide euro and other symbols
\else % if luatex or xetex
  \usepackage{unicode-math} % this also loads fontspec
  \defaultfontfeatures{Scale=MatchLowercase}
  \defaultfontfeatures[\rmfamily]{Ligatures=TeX,Scale=1}
\fi
\usepackage{lmodern}
\ifPDFTeX\else
  % xetex/luatex font selection
\fi
% Use upquote if available, for straight quotes in verbatim environments
\IfFileExists{upquote.sty}{\usepackage{upquote}}{}
\IfFileExists{microtype.sty}{% use microtype if available
  \usepackage[]{microtype}
  \UseMicrotypeSet[protrusion]{basicmath} % disable protrusion for tt fonts
}{}
\makeatletter
\@ifundefined{KOMAClassName}{% if non-KOMA class
  \IfFileExists{parskip.sty}{%
    \usepackage{parskip}
  }{% else
    \setlength{\parindent}{0pt}
    \setlength{\parskip}{6pt plus 2pt minus 1pt}}
}{% if KOMA class
  \KOMAoptions{parskip=half}}
\makeatother
\usepackage{color}
\usepackage{fancyvrb}
\newcommand{\VerbBar}{|}
\newcommand{\VERB}{\Verb[commandchars=\\\{\}]}
\DefineVerbatimEnvironment{Highlighting}{Verbatim}{commandchars=\\\{\}}
% Add ',fontsize=\small' for more characters per line
\usepackage{framed}
\definecolor{shadecolor}{RGB}{248,248,248}
\newenvironment{Shaded}{\begin{snugshade}}{\end{snugshade}}
\newcommand{\AlertTok}[1]{\textcolor[rgb]{0.94,0.16,0.16}{#1}}
\newcommand{\AnnotationTok}[1]{\textcolor[rgb]{0.56,0.35,0.01}{\textbf{\textit{#1}}}}
\newcommand{\AttributeTok}[1]{\textcolor[rgb]{0.13,0.29,0.53}{#1}}
\newcommand{\BaseNTok}[1]{\textcolor[rgb]{0.00,0.00,0.81}{#1}}
\newcommand{\BuiltInTok}[1]{#1}
\newcommand{\CharTok}[1]{\textcolor[rgb]{0.31,0.60,0.02}{#1}}
\newcommand{\CommentTok}[1]{\textcolor[rgb]{0.56,0.35,0.01}{\textit{#1}}}
\newcommand{\CommentVarTok}[1]{\textcolor[rgb]{0.56,0.35,0.01}{\textbf{\textit{#1}}}}
\newcommand{\ConstantTok}[1]{\textcolor[rgb]{0.56,0.35,0.01}{#1}}
\newcommand{\ControlFlowTok}[1]{\textcolor[rgb]{0.13,0.29,0.53}{\textbf{#1}}}
\newcommand{\DataTypeTok}[1]{\textcolor[rgb]{0.13,0.29,0.53}{#1}}
\newcommand{\DecValTok}[1]{\textcolor[rgb]{0.00,0.00,0.81}{#1}}
\newcommand{\DocumentationTok}[1]{\textcolor[rgb]{0.56,0.35,0.01}{\textbf{\textit{#1}}}}
\newcommand{\ErrorTok}[1]{\textcolor[rgb]{0.64,0.00,0.00}{\textbf{#1}}}
\newcommand{\ExtensionTok}[1]{#1}
\newcommand{\FloatTok}[1]{\textcolor[rgb]{0.00,0.00,0.81}{#1}}
\newcommand{\FunctionTok}[1]{\textcolor[rgb]{0.13,0.29,0.53}{\textbf{#1}}}
\newcommand{\ImportTok}[1]{#1}
\newcommand{\InformationTok}[1]{\textcolor[rgb]{0.56,0.35,0.01}{\textbf{\textit{#1}}}}
\newcommand{\KeywordTok}[1]{\textcolor[rgb]{0.13,0.29,0.53}{\textbf{#1}}}
\newcommand{\NormalTok}[1]{#1}
\newcommand{\OperatorTok}[1]{\textcolor[rgb]{0.81,0.36,0.00}{\textbf{#1}}}
\newcommand{\OtherTok}[1]{\textcolor[rgb]{0.56,0.35,0.01}{#1}}
\newcommand{\PreprocessorTok}[1]{\textcolor[rgb]{0.56,0.35,0.01}{\textit{#1}}}
\newcommand{\RegionMarkerTok}[1]{#1}
\newcommand{\SpecialCharTok}[1]{\textcolor[rgb]{0.81,0.36,0.00}{\textbf{#1}}}
\newcommand{\SpecialStringTok}[1]{\textcolor[rgb]{0.31,0.60,0.02}{#1}}
\newcommand{\StringTok}[1]{\textcolor[rgb]{0.31,0.60,0.02}{#1}}
\newcommand{\VariableTok}[1]{\textcolor[rgb]{0.00,0.00,0.00}{#1}}
\newcommand{\VerbatimStringTok}[1]{\textcolor[rgb]{0.31,0.60,0.02}{#1}}
\newcommand{\WarningTok}[1]{\textcolor[rgb]{0.56,0.35,0.01}{\textbf{\textit{#1}}}}
\usepackage{longtable,booktabs,array}
\newcounter{none} % for unnumbered tables
\usepackage{calc} % for calculating minipage widths
% Correct order of tables after \paragraph or \subparagraph
\usepackage{etoolbox}
\makeatletter
\patchcmd\longtable{\par}{\if@noskipsec\mbox{}\fi\par}{}{}
\makeatother
% Allow footnotes in longtable head/foot
\IfFileExists{footnotehyper.sty}{\usepackage{footnotehyper}}{\usepackage{footnote}}
\makesavenoteenv{longtable}
\usepackage{graphicx}
\makeatletter
\newsavebox\pandoc@box
\newcommand*\pandocbounded[1]{% scales image to fit in text height/width
  \sbox\pandoc@box{#1}%
  \Gscale@div\@tempa{\textheight}{\dimexpr\ht\pandoc@box+\dp\pandoc@box\relax}%
  \Gscale@div\@tempb{\linewidth}{\wd\pandoc@box}%
  \ifdim\@tempb\p@<\@tempa\p@\let\@tempa\@tempb\fi% select the smaller of both
  \ifdim\@tempa\p@<\p@\scalebox{\@tempa}{\usebox\pandoc@box}%
  \else\usebox{\pandoc@box}%
  \fi%
}
% Set default figure placement to htbp
\def\fps@figure{htbp}
\makeatother
\setlength{\emergencystretch}{3em} % prevent overfull lines
\providecommand{\tightlist}{%
  \setlength{\itemsep}{0pt}\setlength{\parskip}{0pt}}
\usepackage{booktabs}
\usepackage{longtable}
\usepackage{array}
\usepackage{multirow}
\usepackage{wrapfig}
\usepackage{float}
\usepackage{colortbl}
\usepackage{pdflscape}
\usepackage{tabu}
\usepackage{threeparttable}
\usepackage{threeparttablex}
\usepackage[normalem]{ulem}
\usepackage{makecell}
\usepackage{xcolor}
\usepackage{bookmark}
\IfFileExists{xurl.sty}{\usepackage{xurl}}{} % add URL line breaks if available
\urlstyle{same}
\hypersetup{
  pdftitle={California Housing Price Prediction: Moving Beyond Linearity},
  pdfauthor={Group 1},
  hidelinks,
  pdfcreator={LaTeX via pandoc}}

\title{California Housing Price Prediction: Moving Beyond Linearity}
\usepackage{etoolbox}
\makeatletter
\providecommand{\subtitle}[1]{% add subtitle to \maketitle
  \apptocmd{\@title}{\par {\large #1 \par}}{}{}
}
\makeatother
\subtitle{Nonparametric Regression Project}
\author{Group 1}
\date{2026-05-11}

\begin{document}
\maketitle

{
\setcounter{tocdepth}{3}
\tableofcontents
}
Wesley Gathua, Wendy Oira, Nehama Mburuki, Michelle Andrew, Methuselah
Miningwo, Diana Otieno, Denis Koome, Collins Oduor

\section{Introduction \& Objectives}\label{introduction-objectives}

This project applies the nonparametric and semi-parametric regression
techniques to the California Housing dataset. The central thesis is that
the relationship between housing price and its predictors is
\textbf{inherently nonlinear and spatially interactive}, a claim we test
rigorously by comparing a full ladder of model families.

\textbf{Research Questions}

\begin{enumerate}
\def\labelenumi{\arabic{enumi}.}
\item
  Can ensemble nonlinear models (Random Forest, XGBoost, Neural Network)
  predict house prices more accurately than classical approaches? What
  are the most important nonlinear interactions between location and
  income?
\item
  Which model family --- \textbf{GNPR} (GAM, Kernel/Local Regression,
  Splines, Step Functions, Polynomial) vs.~\textbf{Parametric} (Linear,
  Ridge) vs.~\textbf{Machine Learning} (RF, XGBoost, NN) --- provides
  the best predictive accuracy?
\end{enumerate}

\textbf{Model Ladder}

{\def\LTcaptype{none} % do not increment counter
\begin{longtable}[]{@{}lll@{}}
\toprule\noalign{}
\# & Model & Family \\
\midrule\noalign{}
\endhead
\bottomrule\noalign{}
\endlastfoot
1 & Multiple Linear Regression & Parametric (baseline) \\
2 & Polynomial Regression (degree 3) & GNPR \\
3 & Step Function Regression & GNPR \\
4 & Regression Splines (B-splines, Natural Splines) & GNPR \\
5 & Smoothing Spline (via GAM \texttt{s()}) & GNPR \\
6 & Local Regression (LOESS) & GNPR \\
7 & Generalized Additive Model (GAM + tensor interactions) & GNPR \\
8 & Random Forest & ML Ensemble \\
9 & XGBoost & ML Ensemble \\
10 & Neural Network (MLP) & ML / Deep Learning \\
\end{longtable}
}

\begin{center}\rule{0.5\linewidth}{0.5pt}\end{center}

\section{Package Installation \&
Loading}\label{package-installation-loading}

\begin{Shaded}
\begin{Highlighting}[]
\NormalTok{pkgs }\OtherTok{\textless{}{-}} \FunctionTok{c}\NormalTok{(}
  \StringTok{"tidyverse"}\NormalTok{, }\StringTok{"caret"}\NormalTok{, }\StringTok{"randomForest"}\NormalTok{, }\StringTok{"xgboost"}\NormalTok{, }\StringTok{"mgcv"}\NormalTok{, }\StringTok{"splines"}\NormalTok{, }\StringTok{"np"}\NormalTok{, }\StringTok{"corrplot"}\NormalTok{, }\StringTok{"gridExtra"}\NormalTok{,}
  \StringTok{"Metrics"}\NormalTok{, }\StringTok{"neuralnet"}\NormalTok{, }\StringTok{"ggplot2"}\NormalTok{, }\StringTok{"viridis"}\NormalTok{, }\StringTok{"scales"}\NormalTok{, }\StringTok{"pdp"}\NormalTok{, }\StringTok{"vip"}\NormalTok{, }\StringTok{"spdep"}\NormalTok{, }\StringTok{"sf"}\NormalTok{, }\StringTok{"kableExtra"}\NormalTok{,}
  \StringTok{"GGally"}\NormalTok{, }\StringTok{"patchwork"}\NormalTok{, }\StringTok{"ggrepel"}\NormalTok{, }\StringTok{"patchwork"}\NormalTok{, }\StringTok{"finalfit"}
\NormalTok{)}

\ControlFlowTok{for}\NormalTok{ (pkg }\ControlFlowTok{in}\NormalTok{ pkgs) \{}
  \ControlFlowTok{if}\NormalTok{ (}\SpecialCharTok{!}\NormalTok{(pkg }\SpecialCharTok{\%in\%} \FunctionTok{installed.packages}\NormalTok{())) \{}
    \FunctionTok{install.packages}\NormalTok{(pkg)}
\NormalTok{  \}}
  
  \FunctionTok{library}\NormalTok{(pkg, }\AttributeTok{character.only =}\NormalTok{ T)}
\NormalTok{\}}

\CommentTok{\# Global ggplot theme}
\FunctionTok{theme\_set}\NormalTok{(}\FunctionTok{theme\_minimal}\NormalTok{(}\AttributeTok{base\_size =} \DecValTok{12}\NormalTok{) }\SpecialCharTok{+}
            \FunctionTok{theme}\NormalTok{(}\AttributeTok{plot.title =} \FunctionTok{element\_text}\NormalTok{(}\AttributeTok{face =} \StringTok{"bold"}\NormalTok{, }\AttributeTok{size =} \DecValTok{13}\NormalTok{),}
                  \AttributeTok{plot.subtitle =} \FunctionTok{element\_text}\NormalTok{(}\AttributeTok{color =} \StringTok{"grey40"}\NormalTok{)))}
\end{Highlighting}
\end{Shaded}

\section{Data Loading \& Initial
Inspection}\label{data-loading-initial-inspection}

\begin{Shaded}
\begin{Highlighting}[]
\CommentTok{\# Load dataset (no header row) }
\NormalTok{data }\OtherTok{\textless{}{-}} \FunctionTok{read.csv}\NormalTok{(}\StringTok{"\textasciitilde{}/Downloads/cal\_housing.data"}\NormalTok{, }\AttributeTok{header =} \ConstantTok{FALSE}\NormalTok{)}


\FunctionTok{colnames}\NormalTok{(data) }\OtherTok{\textless{}{-}} \FunctionTok{c}\NormalTok{(}\StringTok{"longitude"}\NormalTok{, }\StringTok{"latitude"}\NormalTok{, }\StringTok{"housing\_median\_age"}\NormalTok{,}
  \StringTok{"total\_rooms"}\NormalTok{, }\StringTok{"total\_bedrooms"}\NormalTok{, }\StringTok{"population"}\NormalTok{,}
  \StringTok{"households"}\NormalTok{, }\StringTok{"median\_income"}\NormalTok{, }\StringTok{"median\_house\_value"}
\NormalTok{)}


\FunctionTok{cat}\NormalTok{(}\StringTok{"Dimensions :"}\NormalTok{, }\FunctionTok{dim}\NormalTok{(data), }\StringTok{"}\SpecialCharTok{\textbackslash{}n}\StringTok{"}\NormalTok{)}
\end{Highlighting}
\end{Shaded}

\begin{verbatim}
#> Dimensions : 20640 9
\end{verbatim}

\begin{Shaded}
\begin{Highlighting}[]
\FunctionTok{cat}\NormalTok{(}\StringTok{"Missing values :"}\NormalTok{, }\FunctionTok{sum}\NormalTok{(}\FunctionTok{is.na}\NormalTok{(data)), }\StringTok{"}\SpecialCharTok{\textbackslash{}n}\StringTok{"}\NormalTok{)}
\end{Highlighting}
\end{Shaded}

\begin{verbatim}
#> Missing values : 0
\end{verbatim}

\begin{Shaded}
\begin{Highlighting}[]
\FunctionTok{head}\NormalTok{(data, }\DecValTok{6}\NormalTok{)}
\end{Highlighting}
\end{Shaded}

\begin{verbatim}
#>   longitude latitude housing_median_age total_rooms total_bedrooms population
#> 1   -122.23    37.88                 41         880            129        322
#> 2   -122.22    37.86                 21        7099           1106       2401
#> 3   -122.24    37.85                 52        1467            190        496
#> 4   -122.25    37.85                 52        1274            235        558
#> 5   -122.25    37.85                 52        1627            280        565
#> 6   -122.25    37.85                 52         919            213        413
#>   households median_income median_house_value
#> 1        126        8.3252             452600
#> 2       1138        8.3014             358500
#> 3        177        7.2574             352100
#> 4        219        5.6431             341300
#> 5        259        3.8462             342200
#> 6        193        4.0368             269700
\end{verbatim}

\begin{Shaded}
\begin{Highlighting}[]
\FunctionTok{summary}\NormalTok{(data)}
\end{Highlighting}
\end{Shaded}

\begin{verbatim}
#>    longitude         latitude     housing_median_age  total_rooms   
#>  Min.   :-124.3   Min.   :32.54   Min.   : 1.00      Min.   :    2  
#>  1st Qu.:-121.8   1st Qu.:33.93   1st Qu.:18.00      1st Qu.: 1448  
#>  Median :-118.5   Median :34.26   Median :29.00      Median : 2127  
#>  Mean   :-119.6   Mean   :35.63   Mean   :28.64      Mean   : 2636  
#>  3rd Qu.:-118.0   3rd Qu.:37.71   3rd Qu.:37.00      3rd Qu.: 3148  
#>  Max.   :-114.3   Max.   :41.95   Max.   :52.00      Max.   :39320  
#>  total_bedrooms     population      households     median_income    
#>  Min.   :   1.0   Min.   :    3   Min.   :   1.0   Min.   : 0.4999  
#>  1st Qu.: 295.0   1st Qu.:  787   1st Qu.: 280.0   1st Qu.: 2.5634  
#>  Median : 435.0   Median : 1166   Median : 409.0   Median : 3.5348  
#>  Mean   : 537.9   Mean   : 1425   Mean   : 499.5   Mean   : 3.8707  
#>  3rd Qu.: 647.0   3rd Qu.: 1725   3rd Qu.: 605.0   3rd Qu.: 4.7432  
#>  Max.   :6445.0   Max.   :35682   Max.   :6082.0   Max.   :15.0001  
#>  median_house_value
#>  Min.   : 14999    
#>  1st Qu.:119600    
#>  Median :179700    
#>  Mean   :206856    
#>  3rd Qu.:264725    
#>  Max.   :500001
\end{verbatim}

\section{Exploratory Data Analysis}\label{exploratory-data-analysis}

\subsection{Distributions of All
Variables}\label{distributions-of-all-variables}

\begin{Shaded}
\begin{Highlighting}[]
\NormalTok{data\_long }\OtherTok{\textless{}{-}}\NormalTok{ data }\SpecialCharTok{\%\textgreater{}\%} 
  \FunctionTok{select}\NormalTok{(}\SpecialCharTok{{-}}\FunctionTok{c}\NormalTok{(longitude,latitude)) }\SpecialCharTok{\%\textgreater{}\%}
  \FunctionTok{pivot\_longer}\NormalTok{(}\FunctionTok{everything}\NormalTok{(), }\AttributeTok{names\_to =} \StringTok{"variable"}\NormalTok{, }\AttributeTok{values\_to =} \StringTok{"value"}\NormalTok{)}

\FunctionTok{ggplot}\NormalTok{(data\_long, }\FunctionTok{aes}\NormalTok{(}\AttributeTok{x =}\NormalTok{ value)) }\SpecialCharTok{+}
  \FunctionTok{geom\_histogram}\NormalTok{(}\AttributeTok{bins =} \DecValTok{40}\NormalTok{, }\AttributeTok{fill =} \StringTok{"steelblue"}\NormalTok{, }\AttributeTok{color =} \StringTok{"white"}\NormalTok{, }\AttributeTok{alpha =} \FloatTok{0.8}\NormalTok{) }\SpecialCharTok{+}
  \CommentTok{\#geom\_vline(xintercept = median(x)) +}
  \FunctionTok{facet\_wrap}\NormalTok{(}\SpecialCharTok{\textasciitilde{}}\NormalTok{variable, }\AttributeTok{scales =} \StringTok{"free"}\NormalTok{) }\SpecialCharTok{+}
  \FunctionTok{theme\_minimal}\NormalTok{() }\SpecialCharTok{+}
  \FunctionTok{labs}\NormalTok{(}\AttributeTok{title =} \StringTok{"Distribution of All Variables"}\NormalTok{)}
\end{Highlighting}
\end{Shaded}

\pandocbounded{\includegraphics[keepaspectratio]{California_Housing_Analysis_files/figure-latex/eda-distributions-1.pdf}}

\begin{itemize}
\tightlist
\item
  The numeric features are skewed to the right apart from
  housing\_median\_age which is bi-modal.
\end{itemize}

\subsection{Correlation Structure}\label{correlation-structure}

\begin{Shaded}
\begin{Highlighting}[]
\NormalTok{cor\_matrix }\OtherTok{\textless{}{-}} \FunctionTok{cor}\NormalTok{(data, }\AttributeTok{use =} \StringTok{"complete.obs"}\NormalTok{)}
\FunctionTok{corrplot}\NormalTok{(}
\NormalTok{  cor\_matrix,}
  \AttributeTok{method      =} \StringTok{"color"}\NormalTok{,}
  \AttributeTok{type        =} \StringTok{"upper"}\NormalTok{,}
  \AttributeTok{tl.cex      =} \FloatTok{0.85}\NormalTok{,}
  \AttributeTok{addCoef.col =} \StringTok{"black"}\NormalTok{,}
  \AttributeTok{number.cex  =} \FloatTok{0.60}\NormalTok{,}
  \AttributeTok{col         =} \FunctionTok{colorRampPalette}\NormalTok{(}\FunctionTok{c}\NormalTok{(}\StringTok{"\#d73027"}\NormalTok{,}\StringTok{"white"}\NormalTok{,}\StringTok{"\#1a6aaa"}\NormalTok{))(}\DecValTok{200}\NormalTok{),}
  \AttributeTok{title       =} \StringTok{"Pearson Correlation Matrix"}\NormalTok{,}
  \AttributeTok{mar         =} \FunctionTok{c}\NormalTok{(}\DecValTok{0}\NormalTok{, }\DecValTok{0}\NormalTok{, }\DecValTok{2}\NormalTok{, }\DecValTok{0}\NormalTok{)}
\NormalTok{)}
\end{Highlighting}
\end{Shaded}

\pandocbounded{\includegraphics[keepaspectratio]{California_Housing_Analysis_files/figure-latex/eda-correlation-1.pdf}}

\begin{itemize}
\tightlist
\item
  There is a very high correlation between total bedrooms, population,
  households and total rooms. Median income has a high correlation with
  median house value.
\end{itemize}

\begin{Shaded}
\begin{Highlighting}[]
\NormalTok{dependent }\OtherTok{\textless{}{-}} \StringTok{"median\_house\_value"}
\NormalTok{explanatory }\OtherTok{\textless{}{-}}\NormalTok{ data }\SpecialCharTok{\%\textgreater{}\%}
  \FunctionTok{select}\NormalTok{(}\SpecialCharTok{{-}}\FunctionTok{c}\NormalTok{(dependent, longitude, latitude)) }\SpecialCharTok{\%\textgreater{}\%}
\NormalTok{  names}

\NormalTok{finalfit}\SpecialCharTok{::}\FunctionTok{ff\_plot}\NormalTok{(data, dependent, explanatory)}
\end{Highlighting}
\end{Shaded}

\pandocbounded{\includegraphics[keepaspectratio]{California_Housing_Analysis_files/figure-latex/unnamed-chunk-1-1.pdf}}

\begin{Shaded}
\begin{Highlighting}[]
\CommentTok{\# Create categorical dependent variable}
\NormalTok{data }\OtherTok{\textless{}{-}}\NormalTok{ data }\SpecialCharTok{\%\textgreater{}\%}
  \FunctionTok{mutate}\NormalTok{(}
    \AttributeTok{house\_value\_group =} \FunctionTok{ntile}\NormalTok{(median\_house\_value, }\DecValTok{4}\NormalTok{) }\SpecialCharTok{\%\textgreater{}\%}
      \FunctionTok{factor}\NormalTok{(}\AttributeTok{labels =} \FunctionTok{c}\NormalTok{(}\StringTok{"Low"}\NormalTok{, }\StringTok{"Medium"}\NormalTok{, }\StringTok{"High"}\NormalTok{, }\StringTok{"Very High"}\NormalTok{))}
\NormalTok{  )}

\NormalTok{dependent }\OtherTok{\textless{}{-}} \StringTok{"house\_value\_group"}

\NormalTok{explanatory }\OtherTok{\textless{}{-}}\NormalTok{ data }\SpecialCharTok{\%\textgreater{}\%}
  \FunctionTok{select}\NormalTok{(}\SpecialCharTok{{-}}\NormalTok{median\_house\_value, }\SpecialCharTok{{-}}\NormalTok{house\_value\_group, }\SpecialCharTok{{-}}\NormalTok{longitude, }\SpecialCharTok{{-}}\NormalTok{latitude) }\SpecialCharTok{\%\textgreater{}\%}
  \FunctionTok{names}\NormalTok{()}

\CommentTok{\# Create finalfit table}
\NormalTok{table1 }\OtherTok{\textless{}{-}}\NormalTok{ data }\SpecialCharTok{\%\textgreater{}\%}
  \FunctionTok{summary\_factorlist}\NormalTok{(}
    \AttributeTok{dependent =}\NormalTok{ dependent,}
    \AttributeTok{explanatory =}\NormalTok{ explanatory,}
    \AttributeTok{p =} \ConstantTok{TRUE}\NormalTok{,}
    \AttributeTok{cont =} \StringTok{"median"}\NormalTok{,}
    \AttributeTok{column =} \ConstantTok{TRUE}\NormalTok{,}
    \AttributeTok{total\_col =} \ConstantTok{TRUE}\NormalTok{,}
    \AttributeTok{na\_include =} \ConstantTok{TRUE}
\NormalTok{  )}

\CommentTok{\# Render as formatted table}
\NormalTok{table1 }\SpecialCharTok{\%\textgreater{}\%}
  \FunctionTok{kbl}\NormalTok{(}
    \AttributeTok{caption =} \StringTok{"Summary Statistics by House Value Group"}\NormalTok{,}
    \AttributeTok{align =} \StringTok{"lcccccc"}\NormalTok{,}
    \AttributeTok{booktabs =} \ConstantTok{FALSE}
\NormalTok{  ) }\SpecialCharTok{\%\textgreater{}\%}
  \FunctionTok{kable\_styling}\NormalTok{(}
    \AttributeTok{bootstrap\_options =} \FunctionTok{c}\NormalTok{(}\StringTok{"striped"}\NormalTok{, }\StringTok{"bordered"}\NormalTok{, }\StringTok{"hover"}\NormalTok{),}
    \AttributeTok{full\_width =} \ConstantTok{FALSE}
\NormalTok{  ) }\SpecialCharTok{\%\textgreater{}\%}
  \FunctionTok{row\_spec}\NormalTok{(}\DecValTok{0}\NormalTok{, }\AttributeTok{bold =} \ConstantTok{TRUE}\NormalTok{)}
\end{Highlighting}
\end{Shaded}

\begin{table}
\centering
\caption{\label{tab:unnamed-chunk-2}Summary Statistics by House Value Group}
\centering
\begin{tabular}[t]{l|c|c|c|c|c|c|l}
\hline
\textbf{label} & \textbf{levels} & \textbf{Low} & \textbf{Medium} & \textbf{High} & \textbf{Very High} & \textbf{Total} & \textbf{p}\\
\hline
housing\_median\_age & Median (IQR) & 28.0 (18.0 to 37.0) & 28.0 (17.0 to 36.0) & 28.0 (19.0 to 36.0) & 31.0 (20.0 to 40.0) & 29.0 (18.0 to 37.0) & <0.001\\
\hline
total\_rooms & Median (IQR) & 1789.5 (1214.8 to 2614.2) & 2020.0 (1391.0 to 3038.2) & 2224.5 (1539.8 to 3299.0) & 2501.0 (1731.8 to 3617.2) & 2127.0 (1447.8 to 3148.0) & <0.001\\
\hline
total\_bedrooms & Median (IQR) & 393.0 (269.0 to 572.0) & 444.0 (301.0 to 661.0) & 458.0 (308.8 to 689.0) & 451.0 (306.0 to 678.0) & 435.0 (295.0 to 647.0) & <0.001\\
\hline
population & Median (IQR) & 1088.0 (727.8 to 1567.0) & 1289.5 (854.8 to 1915.2) & 1227.0 (834.0 to 1810.2) & 1092.0 (752.0 to 1604.0) & 1166.0 (787.0 to 1725.0) & <0.001\\
\hline
households & Median (IQR) & 359.0 (247.0 to 518.0) & 417.0 (285.0 to 614.0) & 431.0 (296.0 to 646.0) & 429.0 (295.0 to 633.0) & 409.0 (280.0 to 605.0) & <0.001\\
\hline
median\_income & Median (IQR) & 2.3 (1.8 to 3.0) & 3.2 (2.6 to 4.0) & 4.0 (3.2 to 4.9) & 5.4 (4.0 to 6.8) & 3.5 (2.6 to 4.7) & <0.001\\
\hline
\end{tabular}
\end{table}

\subsection{Geographic Distribution}\label{geographic-distribution}

\begin{Shaded}
\begin{Highlighting}[]
\FunctionTok{ggplot}\NormalTok{(data, }\FunctionTok{aes}\NormalTok{(}\AttributeTok{x =}\NormalTok{ longitude, }\AttributeTok{y =}\NormalTok{ latitude, }\AttributeTok{color =}\NormalTok{ median\_house\_value }\SpecialCharTok{/} \FloatTok{1e3}\NormalTok{)) }\SpecialCharTok{+}
  \FunctionTok{geom\_point}\NormalTok{(}\AttributeTok{alpha =} \FloatTok{0.25}\NormalTok{, }\AttributeTok{size =} \FloatTok{0.7}\NormalTok{) }\SpecialCharTok{+}
  \FunctionTok{scale\_color\_viridis}\NormalTok{(}\AttributeTok{name =} \StringTok{"Value ($K)"}\NormalTok{, }\AttributeTok{option =} \StringTok{"C"}\NormalTok{, }\AttributeTok{labels =}\NormalTok{ comma) }\SpecialCharTok{+}
  \FunctionTok{coord\_fixed}\NormalTok{(}\AttributeTok{ratio =} \FloatTok{1.3}\NormalTok{) }\SpecialCharTok{+}
  \FunctionTok{labs}\NormalTok{(}
    \AttributeTok{title    =} \StringTok{"Geographic Distribution of Median House Values"}\NormalTok{,}
    \AttributeTok{subtitle =} \StringTok{"Bay Area (\textasciitilde{}37–38°N) and LA (\textasciitilde{}34°N) coastal corridors command the highest prices"}\NormalTok{,}
    \AttributeTok{x =} \StringTok{"Longitude"}\NormalTok{, }\AttributeTok{y =} \StringTok{"Latitude"}
\NormalTok{  )}
\end{Highlighting}
\end{Shaded}

\pandocbounded{\includegraphics[keepaspectratio]{California_Housing_Analysis_files/figure-latex/eda-geo-1.pdf}}

-Expensive neighborhoods tend to be geographically grouped together.
Additionally, expensive neighborhoods are on the western part of
California near the ocean while cheaper houses are located towards the
inland side.Due to urban planning, house prices seems to be regulated to
a maximum of \$500001 and a minimum of approximately \$15000.

\subsection{Spatial Autocorrelation --- Moran's I
Test}\label{spatial-autocorrelation-morans-i-test}

\begin{Shaded}
\begin{Highlighting}[]
\CommentTok{\# Build spatial object and k{-}NN weights (k = 5)}
\NormalTok{spatial\_obj }\OtherTok{\textless{}{-}} \FunctionTok{st\_as\_sf}\NormalTok{(data, }\AttributeTok{coords =} \FunctionTok{c}\NormalTok{(}\StringTok{"longitude"}\NormalTok{, }\StringTok{"latitude"}\NormalTok{), }\AttributeTok{crs =} \DecValTok{4326}\NormalTok{)}
\NormalTok{coords      }\OtherTok{\textless{}{-}} \FunctionTok{st\_coordinates}\NormalTok{(spatial\_obj)}
\NormalTok{knn5        }\OtherTok{\textless{}{-}} \FunctionTok{knn2nb}\NormalTok{(}\FunctionTok{knearneigh}\NormalTok{(coords, }\AttributeTok{k =} \DecValTok{5}\NormalTok{))}
\NormalTok{lw          }\OtherTok{\textless{}{-}} \FunctionTok{nb2listw}\NormalTok{(knn5, }\AttributeTok{style =} \StringTok{"W"}\NormalTok{)}

\CommentTok{\# Moran\textquotesingle{}s I — median income}
\NormalTok{moran\_income }\OtherTok{\textless{}{-}} \FunctionTok{moran.test}\NormalTok{(data}\SpecialCharTok{$}\NormalTok{median\_income, lw)}
\FunctionTok{cat}\NormalTok{(}\StringTok{"Moran\textquotesingle{}s I (median\_income)      :"}\NormalTok{, }\FunctionTok{round}\NormalTok{(moran\_income}\SpecialCharTok{$}\NormalTok{estimate[}\StringTok{"Moran I statistic"}\NormalTok{], }\DecValTok{4}\NormalTok{), }\StringTok{"}\SpecialCharTok{\textbackslash{}n}\StringTok{"}\NormalTok{)}
\end{Highlighting}
\end{Shaded}

\begin{verbatim}
#> Moran's I (median_income)      : 0.5963
\end{verbatim}

\begin{Shaded}
\begin{Highlighting}[]
\FunctionTok{cat}\NormalTok{(}\StringTok{"p{-}value                        :"}\NormalTok{, }\FunctionTok{format}\NormalTok{(moran\_income}\SpecialCharTok{$}\NormalTok{p.value, }\AttributeTok{scientific =} \ConstantTok{TRUE}\NormalTok{), }\StringTok{"}\SpecialCharTok{\textbackslash{}n}\StringTok{"}\NormalTok{)}
\end{Highlighting}
\end{Shaded}

\begin{verbatim}
#> p-value                        : 0e+00
\end{verbatim}

\begin{Shaded}
\begin{Highlighting}[]
\FunctionTok{moran.plot}\NormalTok{(}
\NormalTok{  data}\SpecialCharTok{$}\NormalTok{median\_income, lw,}
  \AttributeTok{labels =} \ConstantTok{FALSE}\NormalTok{,}
  \AttributeTok{main   =} \StringTok{"Moran Scatter Plot — Median Income"}\NormalTok{,}
  \AttributeTok{xlab   =} \StringTok{"Median Income"}\NormalTok{, }\AttributeTok{ylab =} \StringTok{"Spatially Lagged Median Income"}\NormalTok{,}
  \AttributeTok{pch    =} \DecValTok{20}\NormalTok{, }\AttributeTok{col =} \FunctionTok{rgb}\NormalTok{(}\FloatTok{0.18}\NormalTok{, }\FloatTok{0.42}\NormalTok{, }\FloatTok{0.64}\NormalTok{, }\FloatTok{0.3}\NormalTok{)}
\NormalTok{)}
\end{Highlighting}
\end{Shaded}

\pandocbounded{\includegraphics[keepaspectratio]{California_Housing_Analysis_files/figure-latex/moran-income-1.pdf}}

\begin{Shaded}
\begin{Highlighting}[]
\CommentTok{\# Moran\textquotesingle{}s I — house value}
\NormalTok{moran\_value }\OtherTok{\textless{}{-}} \FunctionTok{moran.test}\NormalTok{(data}\SpecialCharTok{$}\NormalTok{median\_house\_value, lw)}
\FunctionTok{cat}\NormalTok{(}\StringTok{"Moran\textquotesingle{}s I (median\_house\_value) :"}\NormalTok{, }\FunctionTok{round}\NormalTok{(moran\_value}\SpecialCharTok{$}\NormalTok{estimate[}\StringTok{"Moran I statistic"}\NormalTok{], }\DecValTok{4}\NormalTok{), }\StringTok{"}\SpecialCharTok{\textbackslash{}n}\StringTok{"}\NormalTok{)}
\end{Highlighting}
\end{Shaded}

\begin{verbatim}
#> Moran's I (median_house_value) : 0.8137
\end{verbatim}

\begin{Shaded}
\begin{Highlighting}[]
\FunctionTok{cat}\NormalTok{(}\StringTok{"p{-}value                        :"}\NormalTok{, }\FunctionTok{format}\NormalTok{(moran\_value}\SpecialCharTok{$}\NormalTok{p.value, }\AttributeTok{scientific =} \ConstantTok{TRUE}\NormalTok{), }\StringTok{"}\SpecialCharTok{\textbackslash{}n}\StringTok{"}\NormalTok{)}
\end{Highlighting}
\end{Shaded}

\begin{verbatim}
#> p-value                        : 0e+00
\end{verbatim}

\begin{Shaded}
\begin{Highlighting}[]
\FunctionTok{moran.plot}\NormalTok{(}
\NormalTok{  data}\SpecialCharTok{$}\NormalTok{median\_house\_value, lw,}
  \AttributeTok{labels =} \ConstantTok{FALSE}\NormalTok{,}
  \AttributeTok{main   =} \StringTok{"Moran Scatter Plot — Median House Value"}\NormalTok{,}
  \AttributeTok{xlab   =} \StringTok{"House Value"}\NormalTok{, }\AttributeTok{ylab =} \StringTok{"Spatially Lagged House Value"}\NormalTok{,}
  \AttributeTok{pch    =} \DecValTok{20}\NormalTok{, }\AttributeTok{col =} \FunctionTok{rgb}\NormalTok{(}\FloatTok{0.77}\NormalTok{, }\FloatTok{0.31}\NormalTok{, }\FloatTok{0.32}\NormalTok{, }\FloatTok{0.3}\NormalTok{)}
\NormalTok{)}
\end{Highlighting}
\end{Shaded}

\pandocbounded{\includegraphics[keepaspectratio]{California_Housing_Analysis_files/figure-latex/moran-value-1.pdf}}

\begin{itemize}
\tightlist
\item
  Both Moran's I statistics are \textbf{highly significant} (p
  \textless{} 2.2e-16), confirming strong positive spatial
  autocorrelation in both income and house values. Nearby block-groups
  tend to have similar values, underpinning why location features
  (latitude, longitude) are critical predictors.
\end{itemize}

\subsection{Nonlinearity: Predictors vs.~House
Value}\label{nonlinearity-predictors-vs.-house-value}

\begin{Shaded}
\begin{Highlighting}[]
\NormalTok{vars }\OtherTok{\textless{}{-}} \FunctionTok{c}\NormalTok{(}
  \StringTok{"longitude"}\NormalTok{, }\StringTok{"latitude"}\NormalTok{, }\StringTok{"housing\_median\_age"}\NormalTok{,}
  \StringTok{"total\_rooms"}\NormalTok{, }\StringTok{"total\_bedrooms"}\NormalTok{, }\StringTok{"population"}\NormalTok{,}
  \StringTok{"households"}\NormalTok{, }\StringTok{"median\_income"}
\NormalTok{)}

\FunctionTok{options}\NormalTok{(}
  \AttributeTok{repr.plot.width =} \DecValTok{18}\NormalTok{,}
  \AttributeTok{repr.plot.height =} \DecValTok{18}
\NormalTok{)}

\CommentTok{\# Create plots}
\NormalTok{plots }\OtherTok{\textless{}{-}} \FunctionTok{lapply}\NormalTok{(vars, }\ControlFlowTok{function}\NormalTok{(v) \{}
  
  \FunctionTok{ggplot}\NormalTok{(data, }\FunctionTok{aes\_string}\NormalTok{(}\AttributeTok{x =}\NormalTok{ v, }\AttributeTok{y =} \StringTok{"median\_house\_value / 1e3"}\NormalTok{)) }\SpecialCharTok{+}
    \FunctionTok{geom\_point}\NormalTok{(}\AttributeTok{alpha =} \FloatTok{0.12}\NormalTok{, }\AttributeTok{color =} \StringTok{"\#2e6da4"}\NormalTok{, }\AttributeTok{size =} \FloatTok{0.7}\NormalTok{) }\SpecialCharTok{+}
    
    \CommentTok{\# LOESS nonlinear fit}
    \FunctionTok{geom\_smooth}\NormalTok{(}
      \AttributeTok{method =} \StringTok{"loess"}\NormalTok{,}
      \AttributeTok{color =} \StringTok{"\#c0392b"}\NormalTok{,}
      \AttributeTok{linewidth =} \FloatTok{1.1}\NormalTok{,}
      \AttributeTok{se =} \ConstantTok{TRUE}
\NormalTok{    ) }\SpecialCharTok{+}
    
    \CommentTok{\# Linear regression fit}
    \FunctionTok{geom\_smooth}\NormalTok{(}
      \AttributeTok{method =} \StringTok{"lm"}\NormalTok{,}
      \AttributeTok{color =} \StringTok{"\#27ae60"}\NormalTok{,}
      \AttributeTok{linewidth =} \FloatTok{0.9}\NormalTok{,}
      \AttributeTok{linetype =} \StringTok{"dashed"}\NormalTok{,}
      \AttributeTok{se =} \ConstantTok{FALSE}
\NormalTok{    ) }\SpecialCharTok{+}
    
    \FunctionTok{labs}\NormalTok{(}
      \AttributeTok{title =} \FunctionTok{paste}\NormalTok{(v, }\StringTok{"vs Median House Value"}\NormalTok{),}
      \AttributeTok{subtitle =} \StringTok{"Red = LOESS | Green dashed = Linear"}\NormalTok{,}
      \AttributeTok{x =}\NormalTok{ v,}
      \AttributeTok{y =} \StringTok{"Median House Value ($K)"}
\NormalTok{    ) }\SpecialCharTok{+}
    
    \FunctionTok{theme\_minimal}\NormalTok{(}\AttributeTok{base\_size =} \DecValTok{11}\NormalTok{)}
\NormalTok{\})}

\CommentTok{\# Combine all plots}
\FunctionTok{wrap\_plots}\NormalTok{(plots, }\AttributeTok{ncol =} \DecValTok{2}\NormalTok{)}
\end{Highlighting}
\end{Shaded}

\pandocbounded{\includegraphics[keepaspectratio]{California_Housing_Analysis_files/figure-latex/unnamed-chunk-3-1.pdf}}

\begin{itemize}
\tightlist
\item
  The LOESS smoothers reveal a \textbf{nonlinear} relationships between
  median house value and several variables. For example, median house
  values increase from low to moderate incomes, then flattens above
  \textasciitilde100K income. This is the primary motivation for using
  polynomial, spline and GAM models rather than linear regression. The
  relationship between median house value and total bedrooms, on the
  other hand, is fairly linear.
\end{itemize}

\section{Feature Engineering \&
Preprocessing}\label{feature-engineering-preprocessing}

\begin{Shaded}
\begin{Highlighting}[]
\CommentTok{\# ── Derived features ───────────────────────────────────────────────────────────}
\NormalTok{data\_clean }\OtherTok{\textless{}{-}}\NormalTok{ data }\SpecialCharTok{\%\textgreater{}\%}
  \FunctionTok{mutate}\NormalTok{(}
    \AttributeTok{rooms\_per\_household      =}\NormalTok{ total\_rooms    }\SpecialCharTok{/}\NormalTok{ households,}
    \AttributeTok{bedrooms\_per\_room        =}\NormalTok{ total\_bedrooms }\SpecialCharTok{/}\NormalTok{ total\_rooms,}
    \AttributeTok{population\_per\_household =}\NormalTok{ population     }\SpecialCharTok{/}\NormalTok{ households}
\NormalTok{  ) }\SpecialCharTok{\%\textgreater{}\%}
  \CommentTok{\# Remove raw aggregate columns that are replaced by per{-}unit ratios}
  \CommentTok{\# (kept for spatial models but removed from feature set to reduce multicollinearity)}
  \FunctionTok{select}\NormalTok{(}\FunctionTok{everything}\NormalTok{()) }\SpecialCharTok{\%\textgreater{}\%}  \CommentTok{\# ensure no old columns}
  \FunctionTok{select}\NormalTok{(longitude, latitude, housing\_median\_age,}
\NormalTok{         total\_rooms, total\_bedrooms, population, households,}
\NormalTok{         median\_income, rooms\_per\_household, bedrooms\_per\_room,}
\NormalTok{         population\_per\_household, median\_house\_value)}

\FunctionTok{cat}\NormalTok{(}\StringTok{"Feature{-}engineered dataset dimensions:"}\NormalTok{, }\FunctionTok{dim}\NormalTok{(data\_clean), }\StringTok{"}\SpecialCharTok{\textbackslash{}n}\StringTok{"}\NormalTok{)}
\end{Highlighting}
\end{Shaded}

\begin{verbatim}
#> Feature-engineered dataset dimensions: 20640 12
\end{verbatim}

\begin{Shaded}
\begin{Highlighting}[]
\FunctionTok{glimpse}\NormalTok{(data\_clean)}
\end{Highlighting}
\end{Shaded}

\begin{verbatim}
#> Rows: 20,640
#> Columns: 12
#> $ longitude                <dbl> -122.23, -122.22, -122.24, -122.25, -122.25, ~
#> $ latitude                 <dbl> 37.88, 37.86, 37.85, 37.85, 37.85, 37.85, 37.~
#> $ housing_median_age       <dbl> 41, 21, 52, 52, 52, 52, 52, 52, 42, 52, 52, 5~
#> $ total_rooms              <dbl> 880, 7099, 1467, 1274, 1627, 919, 2535, 3104,~
#> $ total_bedrooms           <dbl> 129, 1106, 190, 235, 280, 213, 489, 687, 665,~
#> $ population               <dbl> 322, 2401, 496, 558, 565, 413, 1094, 1157, 12~
#> $ households               <dbl> 126, 1138, 177, 219, 259, 193, 514, 647, 595,~
#> $ median_income            <dbl> 8.3252, 8.3014, 7.2574, 5.6431, 3.8462, 4.036~
#> $ rooms_per_household      <dbl> 6.984127, 6.238137, 8.288136, 5.817352, 6.281~
#> $ bedrooms_per_room        <dbl> 0.1465909, 0.1557966, 0.1295160, 0.1844584, 0~
#> $ population_per_household <dbl> 2.555556, 2.109842, 2.802260, 2.547945, 2.181~
#> $ median_house_value       <dbl> 452600, 358500, 352100, 341300, 342200, 26970~
\end{verbatim}

\begin{Shaded}
\begin{Highlighting}[]
\FunctionTok{set.seed}\NormalTok{(}\DecValTok{42}\NormalTok{)}

\CommentTok{\# Feature names used across all models (excludes lat/lon for scaled versions)}
\NormalTok{features\_no\_geo }\OtherTok{\textless{}{-}} \FunctionTok{c}\NormalTok{(}
  \StringTok{"housing\_median\_age"}\NormalTok{, }\StringTok{"median\_income"}\NormalTok{, }\StringTok{"rooms\_per\_household"}\NormalTok{,}
  \StringTok{"bedrooms\_per\_room"}\NormalTok{, }\StringTok{"population\_per\_household"}
\NormalTok{)}

\NormalTok{features\_all }\OtherTok{\textless{}{-}} \FunctionTok{c}\NormalTok{(}\StringTok{"longitude"}\NormalTok{, }\StringTok{"latitude"}\NormalTok{, features\_no\_geo)}

\CommentTok{\# Scale ALL numeric features (for parametric \& NN models) }
\NormalTok{preproc\_all }\OtherTok{\textless{}{-}} \FunctionTok{preProcess}\NormalTok{(}
\NormalTok{  data\_clean[, features\_no\_geo],}
  \AttributeTok{method =} \FunctionTok{c}\NormalTok{(}\StringTok{"center"}\NormalTok{, }\StringTok{"scale"}\NormalTok{)}
\NormalTok{)}
\NormalTok{data\_scaled }\OtherTok{\textless{}{-}} \FunctionTok{predict}\NormalTok{(preproc\_all, data\_clean[, features\_all])}
\NormalTok{data\_scaled}\SpecialCharTok{$}\NormalTok{median\_house\_value }\OtherTok{\textless{}{-}}\NormalTok{ data\_clean}\SpecialCharTok{$}\NormalTok{median\_house\_value}

\CommentTok{\# Train / test split (80 / 20) }
\NormalTok{train\_idx }\OtherTok{\textless{}{-}} \FunctionTok{createDataPartition}\NormalTok{(data\_scaled}\SpecialCharTok{$}\NormalTok{median\_house\_value, }\AttributeTok{p =} \FloatTok{0.8}\NormalTok{, }\AttributeTok{list =} \ConstantTok{FALSE}\NormalTok{)}

\NormalTok{train\_sc  }\OtherTok{\textless{}{-}}\NormalTok{ data\_scaled[ train\_idx, ]}
\NormalTok{test\_sc   }\OtherTok{\textless{}{-}}\NormalTok{ data\_scaled[}\SpecialCharTok{{-}}\NormalTok{train\_idx, ]}

\NormalTok{train\_raw }\OtherTok{\textless{}{-}}\NormalTok{ data\_clean[ train\_idx, ]}
\NormalTok{test\_raw  }\OtherTok{\textless{}{-}}\NormalTok{ data\_clean[}\SpecialCharTok{{-}}\NormalTok{train\_idx, ]}

\NormalTok{X\_train }\OtherTok{\textless{}{-}}\NormalTok{ train\_sc[, features\_all]}
\NormalTok{y\_train }\OtherTok{\textless{}{-}}\NormalTok{ train\_sc}\SpecialCharTok{$}\NormalTok{median\_house\_value}
\NormalTok{X\_test  }\OtherTok{\textless{}{-}}\NormalTok{ test\_sc[,  features\_all]}
\NormalTok{y\_test  }\OtherTok{\textless{}{-}}\NormalTok{ test\_sc}\SpecialCharTok{$}\NormalTok{median\_house\_value}

\FunctionTok{cat}\NormalTok{(}\StringTok{"Training set size :"}\NormalTok{, }\FunctionTok{nrow}\NormalTok{(train\_sc), }\StringTok{"}\SpecialCharTok{\textbackslash{}n}\StringTok{"}\NormalTok{)}
\end{Highlighting}
\end{Shaded}

\begin{verbatim}
#> Training set size : 16513
\end{verbatim}

\begin{Shaded}
\begin{Highlighting}[]
\FunctionTok{cat}\NormalTok{(}\StringTok{"Test set size     :"}\NormalTok{, }\FunctionTok{nrow}\NormalTok{(test\_sc),  }\StringTok{"}\SpecialCharTok{\textbackslash{}n}\StringTok{"}\NormalTok{)}
\end{Highlighting}
\end{Shaded}

\begin{verbatim}
#> Test set size     : 4127
\end{verbatim}

\begin{Shaded}
\begin{Highlighting}[]
\NormalTok{calc\_metrics }\OtherTok{\textless{}{-}} \ControlFlowTok{function}\NormalTok{(actual, predicted, model\_name) \{}
\NormalTok{  rmse\_val }\OtherTok{\textless{}{-}} \FunctionTok{rmse}\NormalTok{(actual, predicted)}
\NormalTok{  mae\_val  }\OtherTok{\textless{}{-}} \FunctionTok{mae}\NormalTok{(actual, predicted)}
\NormalTok{  r2\_val   }\OtherTok{\textless{}{-}} \FunctionTok{cor}\NormalTok{(actual, predicted)}\SpecialCharTok{\^{}}\DecValTok{2}
  \FunctionTok{tibble}\NormalTok{(}
    \AttributeTok{Model =}\NormalTok{ model\_name,}
    \AttributeTok{RMSE  =} \FunctionTok{round}\NormalTok{(rmse\_val, }\DecValTok{2}\NormalTok{),}
    \AttributeTok{MAE   =} \FunctionTok{round}\NormalTok{(mae\_val,  }\DecValTok{2}\NormalTok{),}
    \AttributeTok{R2    =} \FunctionTok{round}\NormalTok{(r2\_val,   }\DecValTok{4}\NormalTok{)}
\NormalTok{  )}
\NormalTok{\}}

\NormalTok{all\_metrics }\OtherTok{\textless{}{-}} \FunctionTok{list}\NormalTok{()}
\NormalTok{all\_preds   }\OtherTok{\textless{}{-}} \FunctionTok{list}\NormalTok{(}\AttributeTok{Actual =}\NormalTok{ y\_test)}
\end{Highlighting}
\end{Shaded}

\section{Model 1 --- Multiple Linear Regression
(Baseline)}\label{model-1-multiple-linear-regression-baseline}

\textbf{Theoretical basis:} The standard parametric model assumes:

\[y_i = \beta_0 + \beta_1 x_{i1} + \cdots + \beta_p x_{ip} + \varepsilon_i\]

This imposes linearity and additivity, a restriction we will
systematically relax in subsequent models.

\begin{Shaded}
\begin{Highlighting}[]
\NormalTok{lm\_formula }\OtherTok{\textless{}{-}} \FunctionTok{as.formula}\NormalTok{(}
  \FunctionTok{paste}\NormalTok{(}\StringTok{"median\_house\_value \textasciitilde{}"}\NormalTok{, }\FunctionTok{paste}\NormalTok{(features\_all, }\AttributeTok{collapse =} \StringTok{" + "}\NormalTok{))}
\NormalTok{)}

\NormalTok{lm\_model }\OtherTok{\textless{}{-}} \FunctionTok{lm}\NormalTok{(lm\_formula, }\AttributeTok{data =}\NormalTok{ train\_sc)}
\FunctionTok{summary}\NormalTok{(lm\_model)}
\end{Highlighting}
\end{Shaded}

\begin{verbatim}
#> 
#> Call:
#> lm(formula = lm_formula, data = train_sc)
#> 
#> Residuals:
#>     Min      1Q  Median      3Q     Max 
#> -595183  -45123  -12626   30988  426459 
#> 
#> Coefficients:
#>                            Estimate Std. Error t value Pr(>|t|)    
#> (Intercept)              -3416759.4    73209.8 -46.671  < 2e-16 ***
#> longitude                  -42544.9      828.8 -51.333  < 2e-16 ***
#> latitude                   -41075.6      791.1 -51.921  < 2e-16 ***
#> housing_median_age          12391.0      587.5  21.092  < 2e-16 ***
#> median_income               83341.4      789.9 105.511  < 2e-16 ***
#> rooms_per_household          9510.9      675.6  14.078  < 2e-16 ***
#> bedrooms_per_room           21806.6      770.6  28.300  < 2e-16 ***
#> population_per_household    -3692.7      497.8  -7.418 1.24e-13 ***
#> ---
#> Signif. codes:  0 '***' 0.001 '**' 0.01 '*' 0.05 '.' 0.1 ' ' 1
#> 
#> Residual standard error: 71350 on 16505 degrees of freedom
#> Multiple R-squared:  0.6146, Adjusted R-squared:  0.6144 
#> F-statistic:  3760 on 7 and 16505 DF,  p-value: < 2.2e-16
\end{verbatim}

\begin{Shaded}
\begin{Highlighting}[]
\FunctionTok{par}\NormalTok{(}\AttributeTok{mfrow =} \FunctionTok{c}\NormalTok{(}\DecValTok{2}\NormalTok{, }\DecValTok{2}\NormalTok{))}
\FunctionTok{plot}\NormalTok{(lm\_model, }\AttributeTok{pch =} \DecValTok{20}\NormalTok{, }\AttributeTok{cex =} \FloatTok{0.4}\NormalTok{, }\AttributeTok{col =} \FunctionTok{rgb}\NormalTok{(}\FloatTok{0.18}\NormalTok{, }\FloatTok{0.42}\NormalTok{, }\FloatTok{0.64}\NormalTok{, }\FloatTok{0.3}\NormalTok{))}
\end{Highlighting}
\end{Shaded}

\pandocbounded{\includegraphics[keepaspectratio]{California_Housing_Analysis_files/figure-latex/lm-diagnostics-1.pdf}}

\begin{Shaded}
\begin{Highlighting}[]
\FunctionTok{par}\NormalTok{(}\AttributeTok{mfrow =} \FunctionTok{c}\NormalTok{(}\DecValTok{1}\NormalTok{, }\DecValTok{1}\NormalTok{))}
\end{Highlighting}
\end{Shaded}

\begin{Shaded}
\begin{Highlighting}[]
\NormalTok{lm\_pred }\OtherTok{\textless{}{-}} \FunctionTok{predict}\NormalTok{(lm\_model, }\AttributeTok{newdata =}\NormalTok{ test\_sc)}

\NormalTok{all\_preds[[}\StringTok{"LinearRegression"}\NormalTok{]] }\OtherTok{\textless{}{-}}\NormalTok{ lm\_pred}
\NormalTok{all\_metrics[[}\StringTok{"LinearRegression"}\NormalTok{]] }\OtherTok{\textless{}{-}} \FunctionTok{calc\_metrics}\NormalTok{(y\_test, lm\_pred, }\StringTok{"1. Linear Regression"}\NormalTok{)}

\FunctionTok{cat}\NormalTok{(}\StringTok{"Linear Regression — RMSE:"}\NormalTok{, all\_metrics}\SpecialCharTok{$}\NormalTok{LinearRegression}\SpecialCharTok{$}\NormalTok{RMSE,}
    \StringTok{"| R²:"}\NormalTok{, all\_metrics}\SpecialCharTok{$}\NormalTok{LinearRegression}\SpecialCharTok{$}\NormalTok{R2, }\StringTok{"}\SpecialCharTok{\textbackslash{}n}\StringTok{"}\NormalTok{)}
\end{Highlighting}
\end{Shaded}

\begin{verbatim}
#> Linear Regression — RMSE: 73005.27 | R²: 0.6134
\end{verbatim}

The residual vs.~fitted plot shows a clear fan-shaped pattern
(heteroscedasticity) and the Q-Q plot reveals heavy tails, both
hallmarks of a violated linearity assumption. The R² of
\textasciitilde0.63 leaves 37\% of variance unexplained.

\begin{center}\rule{0.5\linewidth}{0.5pt}\end{center}

\section{Model 2 --- Polynomial Regression (Degree 3)}\label{polynomial}

Polynomial regression extends the linear model by raising predictors to
powers:

\[y_i = \beta_0 + \beta_1 x_i + \beta_2 x_i^2 + \cdots + \beta_d x_i^d + \varepsilon_i\]

We use \textbf{degree 3} (cubic) for the two strongest nonlinear
predictors (\texttt{median\_income}, \texttt{latitude},
\texttt{longitude}) and linear terms for the remainder, keeping the
model estimable without overfitting.

\begin{Shaded}
\begin{Highlighting}[]
\NormalTok{poly\_model }\OtherTok{\textless{}{-}} \FunctionTok{lm}\NormalTok{(}
\NormalTok{  median\_house\_value }\SpecialCharTok{\textasciitilde{}}
    \FunctionTok{poly}\NormalTok{(median\_income, }\DecValTok{3}\NormalTok{, }\AttributeTok{raw =} \ConstantTok{FALSE}\NormalTok{) }\SpecialCharTok{+}
    \FunctionTok{poly}\NormalTok{(latitude,      }\DecValTok{2}\NormalTok{, }\AttributeTok{raw =} \ConstantTok{FALSE}\NormalTok{) }\SpecialCharTok{+}
    \FunctionTok{poly}\NormalTok{(longitude,     }\DecValTok{4}\NormalTok{, }\AttributeTok{raw =} \ConstantTok{FALSE}\NormalTok{) }\SpecialCharTok{+}
\NormalTok{    housing\_median\_age }\SpecialCharTok{+}
\NormalTok{    rooms\_per\_household }\SpecialCharTok{+}\NormalTok{ bedrooms\_per\_room }\SpecialCharTok{+}
\NormalTok{    population\_per\_household,}
  \AttributeTok{data =}\NormalTok{ train\_sc}
\NormalTok{)}

\FunctionTok{summary}\NormalTok{(poly\_model)}
\end{Highlighting}
\end{Shaded}

\begin{verbatim}
#> 
#> Call:
#> lm(formula = median_house_value ~ poly(median_income, 3, raw = FALSE) + 
#>     poly(latitude, 2, raw = FALSE) + poly(longitude, 4, raw = FALSE) + 
#>     housing_median_age + rooms_per_household + bedrooms_per_room + 
#>     population_per_household, data = train_sc)
#> 
#> Residuals:
#>     Min      1Q  Median      3Q     Max 
#> -397883  -44177  -10870   30649  446455 
#> 
#> Coefficients:
#>                                        Estimate Std. Error t value Pr(>|t|)    
#> (Intercept)                           2.065e+05  5.423e+02 380.703  < 2e-16 ***
#> poly(median_income, 3, raw = FALSE)1  1.105e+07  1.024e+05 107.883  < 2e-16 ***
#> poly(median_income, 3, raw = FALSE)2 -1.464e+06  7.548e+04 -19.392  < 2e-16 ***
#> poly(median_income, 3, raw = FALSE)3 -1.207e+06  7.014e+04 -17.209  < 2e-16 ***
#> poly(latitude, 2, raw = FALSE)1      -1.144e+07  2.498e+05 -45.808  < 2e-16 ***
#> poly(latitude, 2, raw = FALSE)2       9.692e+05  9.484e+04  10.220  < 2e-16 ***
#> poly(longitude, 4, raw = FALSE)1     -1.125e+07  2.466e+05 -45.636  < 2e-16 ***
#> poly(longitude, 4, raw = FALSE)2     -5.094e+05  8.285e+04  -6.149 7.98e-10 ***
#> poly(longitude, 4, raw = FALSE)3      2.937e+05  7.928e+04   3.704 0.000213 ***
#> poly(longitude, 4, raw = FALSE)4      2.207e+04  7.302e+04   0.302 0.762509    
#> housing_median_age                    1.449e+04  5.836e+02  24.833  < 2e-16 ***
#> rooms_per_household                   9.454e+03  6.616e+02  14.289  < 2e-16 ***
#> bedrooms_per_room                     2.562e+04  8.021e+02  31.939  < 2e-16 ***
#> population_per_household             -3.873e+03  4.863e+02  -7.964 1.77e-15 ***
#> ---
#> Signif. codes:  0 '***' 0.001 '**' 0.01 '*' 0.05 '.' 0.1 ' ' 1
#> 
#> Residual standard error: 69690 on 16499 degrees of freedom
#> Multiple R-squared:  0.6324, Adjusted R-squared:  0.6321 
#> F-statistic:  2183 on 13 and 16499 DF,  p-value: < 2.2e-16
\end{verbatim}

\begin{Shaded}
\begin{Highlighting}[]
\CommentTok{\# Visualise cubic polynomial fit for median\_income (holding other vars at median)}
\NormalTok{income\_grid }\OtherTok{\textless{}{-}} \FunctionTok{data.frame}\NormalTok{(}
  \AttributeTok{median\_income            =} \FunctionTok{seq}\NormalTok{(}\FunctionTok{min}\NormalTok{(train\_sc}\SpecialCharTok{$}\NormalTok{median\_income), }\FunctionTok{max}\NormalTok{(train\_sc}\SpecialCharTok{$}\NormalTok{median\_income), }\AttributeTok{length =} \DecValTok{200}\NormalTok{),}
  \AttributeTok{latitude                 =} \FunctionTok{median}\NormalTok{(train\_sc}\SpecialCharTok{$}\NormalTok{latitude),}
  \AttributeTok{longitude                =} \FunctionTok{median}\NormalTok{(train\_sc}\SpecialCharTok{$}\NormalTok{longitude),}
  \AttributeTok{housing\_median\_age       =} \FunctionTok{median}\NormalTok{(train\_sc}\SpecialCharTok{$}\NormalTok{housing\_median\_age),}
  \AttributeTok{rooms\_per\_household      =} \FunctionTok{median}\NormalTok{(train\_sc}\SpecialCharTok{$}\NormalTok{rooms\_per\_household),}
  \AttributeTok{bedrooms\_per\_room        =} \FunctionTok{median}\NormalTok{(train\_sc}\SpecialCharTok{$}\NormalTok{bedrooms\_per\_room),}
  \AttributeTok{population\_per\_household =} \FunctionTok{median}\NormalTok{(train\_sc}\SpecialCharTok{$}\NormalTok{population\_per\_household)}
\NormalTok{)}

\NormalTok{income\_pred }\OtherTok{\textless{}{-}} \FunctionTok{predict}\NormalTok{(poly\_model, }\AttributeTok{newdata =}\NormalTok{ income\_grid, }\AttributeTok{interval =} \StringTok{"confidence"}\NormalTok{)}
\NormalTok{income\_grid}\SpecialCharTok{$}\NormalTok{fit  }\OtherTok{\textless{}{-}}\NormalTok{ income\_pred[, }\StringTok{"fit"}\NormalTok{]}
\NormalTok{income\_grid}\SpecialCharTok{$}\NormalTok{lwr  }\OtherTok{\textless{}{-}}\NormalTok{ income\_pred[, }\StringTok{"lwr"}\NormalTok{]}
\NormalTok{income\_grid}\SpecialCharTok{$}\NormalTok{upr  }\OtherTok{\textless{}{-}}\NormalTok{ income\_pred[, }\StringTok{"upr"}\NormalTok{]}

\CommentTok{\# Back{-}transform income from scaled to original for labelling}
\NormalTok{income\_orig }\OtherTok{\textless{}{-}}\NormalTok{ income\_grid}\SpecialCharTok{$}\NormalTok{median\_income }\SpecialCharTok{*}
  \FunctionTok{attr}\NormalTok{(preproc\_all}\SpecialCharTok{$}\NormalTok{mean, }\StringTok{"scaled:scale"}\NormalTok{)[}\StringTok{"median\_income"}\NormalTok{] }\SpecialCharTok{+}
\NormalTok{  preproc\_all}\SpecialCharTok{$}\NormalTok{mean[}\StringTok{"median\_income"}\NormalTok{]}

\FunctionTok{ggplot}\NormalTok{(income\_grid, }\FunctionTok{aes}\NormalTok{(}\AttributeTok{x =}\NormalTok{ income\_grid}\SpecialCharTok{$}\NormalTok{median\_income)) }\SpecialCharTok{+}
  \FunctionTok{geom\_ribbon}\NormalTok{(}\FunctionTok{aes}\NormalTok{(}\AttributeTok{ymin =}\NormalTok{ lwr, }\AttributeTok{ymax =}\NormalTok{ upr), }\AttributeTok{fill =} \StringTok{"\#2e6da4"}\NormalTok{, }\AttributeTok{alpha =} \FloatTok{0.2}\NormalTok{) }\SpecialCharTok{+}
  \FunctionTok{geom\_line}\NormalTok{(}\FunctionTok{aes}\NormalTok{(}\AttributeTok{y =}\NormalTok{ fit), }\AttributeTok{color =} \StringTok{"\#2e6da4"}\NormalTok{, }\AttributeTok{linewidth =} \FloatTok{1.2}\NormalTok{) }\SpecialCharTok{+}
  \FunctionTok{labs}\NormalTok{(}
    \AttributeTok{title    =} \StringTok{"Polynomial Regression (Degree 3): Marginal Effect of Median Income"}\NormalTok{,}
    \AttributeTok{subtitle =} \StringTok{"Other predictors held at their medians. Shaded band = 95\% CI"}\NormalTok{,}
    \AttributeTok{x =} \StringTok{"Median Income (scaled)"}\NormalTok{, }\AttributeTok{y =} \StringTok{"Fitted House Value ($)"}
\NormalTok{  )}
\end{Highlighting}
\end{Shaded}

\pandocbounded{\includegraphics[keepaspectratio]{California_Housing_Analysis_files/figure-latex/poly-plot-1.pdf}}

\begin{Shaded}
\begin{Highlighting}[]
\NormalTok{poly\_pred }\OtherTok{\textless{}{-}} \FunctionTok{predict}\NormalTok{(poly\_model, }\AttributeTok{newdata =}\NormalTok{ test\_sc)}

\NormalTok{all\_preds[[}\StringTok{"PolynomialRegression"}\NormalTok{]] }\OtherTok{\textless{}{-}}\NormalTok{ poly\_pred}
\NormalTok{all\_metrics[[}\StringTok{"PolynomialRegression"}\NormalTok{]] }\OtherTok{\textless{}{-}} \FunctionTok{calc\_metrics}\NormalTok{(y\_test, poly\_pred, }\StringTok{"2. Polynomial Regression (d=3)"}\NormalTok{)}

\FunctionTok{cat}\NormalTok{(}\StringTok{"Polynomial Regression — RMSE:"}\NormalTok{, all\_metrics}\SpecialCharTok{$}\NormalTok{PolynomialRegression}\SpecialCharTok{$}\NormalTok{RMSE,}
    \StringTok{"| R²:"}\NormalTok{, all\_metrics}\SpecialCharTok{$}\NormalTok{PolynomialRegression}\SpecialCharTok{$}\NormalTok{R2, }\StringTok{"}\SpecialCharTok{\textbackslash{}n}\StringTok{"}\NormalTok{)}
\end{Highlighting}
\end{Shaded}

\begin{verbatim}
#> Polynomial Regression — RMSE: 70847.09 | R²: 0.6357
\end{verbatim}

Polynomial regression improves over the linear baseline by capturing the
concave income-price relationship. However, polynomials impose a
\textbf{global} structure,that is, they behave erratically near the
boundary of the predictor space, which motivates local methods like
splines.

\begin{center}\rule{0.5\linewidth}{0.5pt}\end{center}

\section{Model 3 --- Step Function Regression}\label{step}

Step functions partition the range of a predictor into \(K\) bins,
fitting a constant in each:

\[y_i = \beta_0 + \beta_1 C_1(x_i) + \beta_2 C_2(x_i) + \cdots + \beta_K C_K(x_i) + \varepsilon_i\]

where \(C_k(x) = \mathbf{1}(c_{k-1} \le x < c_k)\). Cuts are placed at
\textbf{quartiles} of each predictor.

\begin{Shaded}
\begin{Highlighting}[]
\CommentTok{\# Create step (piecewise{-}constant) predictors using quartile cut{-}points}
\NormalTok{train\_step }\OtherTok{\textless{}{-}}\NormalTok{ train\_sc }\SpecialCharTok{\%\textgreater{}\%}
  \FunctionTok{mutate}\NormalTok{(}
    \AttributeTok{income\_step =} \FunctionTok{cut}\NormalTok{(median\_income, }\AttributeTok{breaks =} \FunctionTok{quantile}\NormalTok{(median\_income, }\AttributeTok{probs =} \DecValTok{0}\SpecialCharTok{:}\DecValTok{4}\SpecialCharTok{/}\DecValTok{4}\NormalTok{),}
                      \AttributeTok{include.lowest =} \ConstantTok{TRUE}\NormalTok{),}
    \AttributeTok{lat\_step    =} \FunctionTok{cut}\NormalTok{(latitude,      }\AttributeTok{breaks =} \FunctionTok{quantile}\NormalTok{(latitude,      }\AttributeTok{probs =} \DecValTok{0}\SpecialCharTok{:}\DecValTok{4}\SpecialCharTok{/}\DecValTok{4}\NormalTok{),}
                      \AttributeTok{include.lowest =} \ConstantTok{TRUE}\NormalTok{),}
    \AttributeTok{lon\_step    =} \FunctionTok{cut}\NormalTok{(longitude,     }\AttributeTok{breaks =} \FunctionTok{quantile}\NormalTok{(longitude,     }\AttributeTok{probs =} \DecValTok{0}\SpecialCharTok{:}\DecValTok{4}\SpecialCharTok{/}\DecValTok{4}\NormalTok{),}
                      \AttributeTok{include.lowest =} \ConstantTok{TRUE}\NormalTok{),}
    \AttributeTok{age\_step    =} \FunctionTok{cut}\NormalTok{(housing\_median\_age, }\AttributeTok{breaks =} \FunctionTok{quantile}\NormalTok{(housing\_median\_age, }\AttributeTok{probs =} \DecValTok{0}\SpecialCharTok{:}\DecValTok{3}\SpecialCharTok{/}\DecValTok{3}\NormalTok{),}
                      \AttributeTok{include.lowest =} \ConstantTok{TRUE}\NormalTok{)}
\NormalTok{  )}

\CommentTok{\# Save cut{-}points for test set}
\NormalTok{income\_breaks }\OtherTok{\textless{}{-}} \FunctionTok{quantile}\NormalTok{(train\_sc}\SpecialCharTok{$}\NormalTok{median\_income,      }\AttributeTok{probs =} \DecValTok{0}\SpecialCharTok{:}\DecValTok{4}\SpecialCharTok{/}\DecValTok{4}\NormalTok{)}
\NormalTok{lat\_breaks    }\OtherTok{\textless{}{-}} \FunctionTok{quantile}\NormalTok{(train\_sc}\SpecialCharTok{$}\NormalTok{latitude,            }\AttributeTok{probs =} \DecValTok{0}\SpecialCharTok{:}\DecValTok{4}\SpecialCharTok{/}\DecValTok{4}\NormalTok{)}
\NormalTok{lon\_breaks    }\OtherTok{\textless{}{-}} \FunctionTok{quantile}\NormalTok{(train\_sc}\SpecialCharTok{$}\NormalTok{longitude,           }\AttributeTok{probs =} \DecValTok{0}\SpecialCharTok{:}\DecValTok{4}\SpecialCharTok{/}\DecValTok{4}\NormalTok{)}
\NormalTok{age\_breaks    }\OtherTok{\textless{}{-}} \FunctionTok{quantile}\NormalTok{(train\_sc}\SpecialCharTok{$}\NormalTok{housing\_median\_age,  }\AttributeTok{probs =} \DecValTok{0}\SpecialCharTok{:}\DecValTok{3}\SpecialCharTok{/}\DecValTok{3}\NormalTok{)}

\NormalTok{test\_step }\OtherTok{\textless{}{-}}\NormalTok{ test\_sc }\SpecialCharTok{\%\textgreater{}\%}
  \FunctionTok{mutate}\NormalTok{(}
    \AttributeTok{income\_step =} \FunctionTok{cut}\NormalTok{(median\_income,      }\AttributeTok{breaks =}\NormalTok{ income\_breaks, }\AttributeTok{include.lowest =} \ConstantTok{TRUE}\NormalTok{),}
    \AttributeTok{lat\_step    =} \FunctionTok{cut}\NormalTok{(latitude,           }\AttributeTok{breaks =}\NormalTok{ lat\_breaks,    }\AttributeTok{include.lowest =} \ConstantTok{TRUE}\NormalTok{),}
    \AttributeTok{lon\_step    =} \FunctionTok{cut}\NormalTok{(longitude,          }\AttributeTok{breaks =}\NormalTok{ lon\_breaks,    }\AttributeTok{include.lowest =} \ConstantTok{TRUE}\NormalTok{),}
    \AttributeTok{age\_step    =} \FunctionTok{cut}\NormalTok{(housing\_median\_age, }\AttributeTok{breaks =}\NormalTok{ age\_breaks,    }\AttributeTok{include.lowest =} \ConstantTok{TRUE}\NormalTok{)}
\NormalTok{  )}

\NormalTok{step\_model }\OtherTok{\textless{}{-}} \FunctionTok{lm}\NormalTok{(}
\NormalTok{  median\_house\_value }\SpecialCharTok{\textasciitilde{}}
\NormalTok{    income\_step }\SpecialCharTok{+}\NormalTok{ lat\_step }\SpecialCharTok{+}\NormalTok{ lon\_step }\SpecialCharTok{+}\NormalTok{ age\_step }\SpecialCharTok{+}
\NormalTok{    rooms\_per\_household }\SpecialCharTok{+}\NormalTok{ bedrooms\_per\_room }\SpecialCharTok{+}\NormalTok{ population\_per\_household,}
  \AttributeTok{data =}\NormalTok{ train\_step}
\NormalTok{)}

\FunctionTok{summary}\NormalTok{(step\_model)}
\end{Highlighting}
\end{Shaded}

\begin{verbatim}
#> 
#> Call:
#> lm(formula = median_house_value ~ income_step + lat_step + lon_step + 
#>     age_step + rooms_per_household + bedrooms_per_room + population_per_household, 
#>     data = train_step)
#> 
#> Residuals:
#>     Min      1Q  Median      3Q     Max 
#> -307135  -53114  -13754   36887  403974 
#> 
#> Coefficients:
#>                             Estimate Std. Error t value Pr(>|t|)    
#> (Intercept)                 224611.3     4538.2  49.494  < 2e-16 ***
#> income_step(-0.688,-0.181]   52688.5     1882.7  27.986  < 2e-16 ***
#> income_step(-0.181,0.457]    97721.0     2064.6  47.333  < 2e-16 ***
#> income_step(0.457,5.86]     196466.7     2454.3  80.051  < 2e-16 ***
#> lat_step(33.9,34.2]          -1404.3     1997.6  -0.703    0.482    
#> lat_step(34.2,37.7]         -77529.9     3463.6 -22.384  < 2e-16 ***
#> lat_step(37.7,42]          -103675.8     3837.5 -27.017  < 2e-16 ***
#> lon_step(-121.8,-118.5]     -68030.4     1981.7 -34.329  < 2e-16 ***
#> lon_step(-118.5,-118]       -82798.3     3552.3 -23.308  < 2e-16 ***
#> lon_step(-118,-114.3]      -129619.5     3790.1 -34.199  < 2e-16 ***
#> age_step(-0.528,0.505]        9709.6     1577.4   6.156 7.65e-10 ***
#> age_step(0.505,1.86]         23807.4     1750.4  13.601  < 2e-16 ***
#> rooms_per_household           8559.4      755.3  11.332  < 2e-16 ***
#> bedrooms_per_room            12819.1      912.2  14.054  < 2e-16 ***
#> population_per_household     -2951.1      569.9  -5.179 2.26e-07 ***
#> ---
#> Signif. codes:  0 '***' 0.001 '**' 0.01 '*' 0.05 '.' 0.1 ' ' 1
#> 
#> Residual standard error: 81720 on 16498 degrees of freedom
#> Multiple R-squared:  0.4945, Adjusted R-squared:  0.4941 
#> F-statistic:  1153 on 14 and 16498 DF,  p-value: < 2.2e-16
\end{verbatim}

\begin{Shaded}
\begin{Highlighting}[]
\CommentTok{\# Visualise step fit for median income}
\NormalTok{step\_income\_df }\OtherTok{\textless{}{-}}\NormalTok{ train\_step }\SpecialCharTok{\%\textgreater{}\%}
  \FunctionTok{group\_by}\NormalTok{(income\_step) }\SpecialCharTok{\%\textgreater{}\%}
  \FunctionTok{summarise}\NormalTok{(}
    \AttributeTok{mean\_income  =} \FunctionTok{mean}\NormalTok{(median\_income),}
    \AttributeTok{mean\_value   =} \FunctionTok{mean}\NormalTok{(median\_house\_value),}
    \AttributeTok{fitted\_value =} \FunctionTok{mean}\NormalTok{(}\FunctionTok{fitted}\NormalTok{(step\_model)[}\FunctionTok{cur\_group\_rows}\NormalTok{()])}
\NormalTok{  )}

\FunctionTok{ggplot}\NormalTok{(train\_step, }\FunctionTok{aes}\NormalTok{(}\AttributeTok{x =}\NormalTok{ median\_income, }\AttributeTok{y =}\NormalTok{ median\_house\_value)) }\SpecialCharTok{+}
  \FunctionTok{geom\_point}\NormalTok{(}\AttributeTok{alpha =} \FloatTok{0.05}\NormalTok{, }\AttributeTok{color =} \StringTok{"\#2e6da4"}\NormalTok{, }\AttributeTok{size =} \FloatTok{0.5}\NormalTok{) }\SpecialCharTok{+}
  \FunctionTok{geom\_step}\NormalTok{(}\AttributeTok{data =}\NormalTok{ step\_income\_df, }\FunctionTok{aes}\NormalTok{(}\AttributeTok{x =}\NormalTok{ mean\_income, }\AttributeTok{y =}\NormalTok{ fitted\_value),}
            \AttributeTok{color =} \StringTok{"\#c0392b"}\NormalTok{, }\AttributeTok{linewidth =} \FloatTok{1.3}\NormalTok{, }\AttributeTok{direction =} \StringTok{"mid"}\NormalTok{) }\SpecialCharTok{+}
  \FunctionTok{labs}\NormalTok{(}
    \AttributeTok{title    =} \StringTok{"Step Function Regression: Piecewise{-}Constant Fit for Median Income"}\NormalTok{,}
    \AttributeTok{subtitle =} \StringTok{"4 income bins (quartile cut{-}points). Red step line = model fitted values per bin"}\NormalTok{,}
    \AttributeTok{x =} \StringTok{"Median Income (scaled)"}\NormalTok{, }\AttributeTok{y =} \StringTok{"House Value ($)"}
\NormalTok{  )}
\end{Highlighting}
\end{Shaded}

\pandocbounded{\includegraphics[keepaspectratio]{California_Housing_Analysis_files/figure-latex/step-plot-1.pdf}}

\begin{Shaded}
\begin{Highlighting}[]
\NormalTok{step\_pred }\OtherTok{\textless{}{-}} \FunctionTok{predict}\NormalTok{(step\_model, }\AttributeTok{newdata =}\NormalTok{ test\_step)}

\NormalTok{all\_preds[[}\StringTok{"StepFunction"}\NormalTok{]] }\OtherTok{\textless{}{-}}\NormalTok{ step\_pred}
\NormalTok{all\_metrics[[}\StringTok{"StepFunction"}\NormalTok{]] }\OtherTok{\textless{}{-}} \FunctionTok{calc\_metrics}\NormalTok{(y\_test, step\_pred, }\StringTok{"3. Step Function"}\NormalTok{)}

\FunctionTok{cat}\NormalTok{(}\StringTok{"Step Function — RMSE:"}\NormalTok{, all\_metrics}\SpecialCharTok{$}\NormalTok{StepFunction}\SpecialCharTok{$}\NormalTok{RMSE,}
    \StringTok{"| R²:"}\NormalTok{, all\_metrics}\SpecialCharTok{$}\NormalTok{StepFunction}\SpecialCharTok{$}\NormalTok{R2, }\StringTok{"}\SpecialCharTok{\textbackslash{}n}\StringTok{"}\NormalTok{)}
\end{Highlighting}
\end{Shaded}

\begin{verbatim}
#> Step Function — RMSE: 83018.9 | R²: 0.4998
\end{verbatim}

Step functions are easy to interpret --- each coefficient represents the
average house value increment in that bin relative to the reference.
However, the discontinuities at knot boundaries are unrealistic; house
prices do not jump discretely. Splines address this by enforcing
continuity.

\begin{center}\rule{0.5\linewidth}{0.5pt}\end{center}

\section{Model 4 --- Generalized Additive Model (GAM)}\label{gam}

GAMs replace linear predictors with flexible smooth functions while
maintaining additivity:

\[y_i = \beta_0 + \sum_{j=1}^{p} f_j(x_{ij}) + \varepsilon_i\]

We extend this with \textbf{tensor product interaction terms}
\texttt{ti()} to capture the income × location interaction that an
additive model would otherwise miss.

\begin{Shaded}
\begin{Highlighting}[]
\NormalTok{gam\_model }\OtherTok{\textless{}{-}}\NormalTok{ mgcv}\SpecialCharTok{::}\FunctionTok{gam}\NormalTok{(}
\NormalTok{  median\_house\_value }\SpecialCharTok{\textasciitilde{}}
    \FunctionTok{s}\NormalTok{(median\_income) }\SpecialCharTok{+}
    \FunctionTok{s}\NormalTok{(latitude) }\SpecialCharTok{+}
    \FunctionTok{s}\NormalTok{(longitude) }\SpecialCharTok{+}
    \FunctionTok{s}\NormalTok{(housing\_median\_age) }\SpecialCharTok{+}
    \FunctionTok{s}\NormalTok{(rooms\_per\_household) }\SpecialCharTok{+}
    \FunctionTok{s}\NormalTok{(population\_per\_household) }\SpecialCharTok{+}
    \FunctionTok{s}\NormalTok{(bedrooms\_per\_room) }\SpecialCharTok{+}
    \FunctionTok{ti}\NormalTok{(latitude, longitude) }\SpecialCharTok{+}           \CommentTok{\# spatial interaction term}
    \FunctionTok{ti}\NormalTok{(median\_income, latitude),         }\CommentTok{\# income × location interaction}
  \AttributeTok{data   =}\NormalTok{ train\_sc,}
  \AttributeTok{method =} \StringTok{"REML"}
\NormalTok{)}

\FunctionTok{summary}\NormalTok{(gam\_model)}
\end{Highlighting}
\end{Shaded}

\begin{verbatim}
#> 
#> Family: gaussian 
#> Link function: identity 
#> 
#> Formula:
#> median_house_value ~ s(median_income) + s(latitude) + s(longitude) + 
#>     s(housing_median_age) + s(rooms_per_household) + s(population_per_household) + 
#>     s(bedrooms_per_room) + ti(latitude, longitude) + ti(median_income, 
#>     latitude)
#> 
#> Parametric coefficients:
#>             Estimate Std. Error t value Pr(>|t|)    
#> (Intercept)   198380      34846   5.693 1.27e-08 ***
#> ---
#> Signif. codes:  0 '***' 0.001 '**' 0.01 '*' 0.05 '.' 0.1 ' ' 1
#> 
#> Approximate significance of smooth terms:
#>                                edf Ref.df       F p-value    
#> s(median_income)             8.048  8.756 431.200  <2e-16 ***
#> s(latitude)                  8.464  8.797  37.580  <2e-16 ***
#> s(longitude)                 8.525  8.732  59.972  <2e-16 ***
#> s(housing_median_age)        8.367  8.872  26.143  <2e-16 ***
#> s(rooms_per_household)       7.844  8.568  30.375  <2e-16 ***
#> s(population_per_household)  8.661  8.960 398.085  <2e-16 ***
#> s(bedrooms_per_room)         5.344  6.590   6.875  <2e-16 ***
#> ti(latitude,longitude)      12.818 13.281  81.965  <2e-16 ***
#> ti(median_income,latitude)   7.733  9.352   7.077  <2e-16 ***
#> ---
#> Signif. codes:  0 '***' 0.001 '**' 0.01 '*' 0.05 '.' 0.1 ' ' 1
#> 
#> R-sq.(adj) =  0.743   Deviance explained = 74.4%
#> -REML = 2.0469e+05  Scale est. = 3.3916e+09  n = 16513
\end{verbatim}

\begin{Shaded}
\begin{Highlighting}[]
\FunctionTok{par}\NormalTok{(}\AttributeTok{mfrow =} \FunctionTok{c}\NormalTok{(}\DecValTok{3}\NormalTok{, }\DecValTok{3}\NormalTok{))}
\FunctionTok{plot}\NormalTok{(}
\NormalTok{  gam\_model, }\AttributeTok{pages =} \DecValTok{0}\NormalTok{,}
  \AttributeTok{shade     =} \ConstantTok{TRUE}\NormalTok{,}
  \AttributeTok{col       =} \StringTok{"\#2e6da4"}\NormalTok{,}
  \AttributeTok{shade.col =} \StringTok{"lightblue"}\NormalTok{,}
  \AttributeTok{seWithMean =} \ConstantTok{TRUE}\NormalTok{,}
  \AttributeTok{residuals  =} \ConstantTok{TRUE}\NormalTok{,}
  \AttributeTok{pch =} \DecValTok{20}\NormalTok{, }\AttributeTok{cex =} \FloatTok{0.15}
\NormalTok{)}
\end{Highlighting}
\end{Shaded}

\pandocbounded{\includegraphics[keepaspectratio]{California_Housing_Analysis_files/figure-latex/gam-smooth-plots-1.pdf}}

\begin{Shaded}
\begin{Highlighting}[]
\FunctionTok{par}\NormalTok{(}\AttributeTok{mfrow =} \FunctionTok{c}\NormalTok{(}\DecValTok{1}\NormalTok{, }\DecValTok{1}\NormalTok{))}
\end{Highlighting}
\end{Shaded}

\begin{Shaded}
\begin{Highlighting}[]
\NormalTok{gam\_pred }\OtherTok{\textless{}{-}} \FunctionTok{predict}\NormalTok{(gam\_model, }\AttributeTok{newdata =}\NormalTok{ test\_sc)}

\NormalTok{all\_preds[[}\StringTok{"GAM"}\NormalTok{]] }\OtherTok{\textless{}{-}}\NormalTok{ gam\_pred}
\NormalTok{all\_metrics[[}\StringTok{"GAM"}\NormalTok{]] }\OtherTok{\textless{}{-}} \FunctionTok{calc\_metrics}\NormalTok{(y\_test, gam\_pred, }\StringTok{"7. GAM (+ tensor interactions)"}\NormalTok{)}

\FunctionTok{cat}\NormalTok{(}\StringTok{"GAM — RMSE:"}\NormalTok{, all\_metrics}\SpecialCharTok{$}\NormalTok{GAM}\SpecialCharTok{$}\NormalTok{RMSE, }\StringTok{"| R²:"}\NormalTok{, all\_metrics}\SpecialCharTok{$}\NormalTok{GAM}\SpecialCharTok{$}\NormalTok{R2, }\StringTok{"}\SpecialCharTok{\textbackslash{}n}\StringTok{"}\NormalTok{)}
\end{Highlighting}
\end{Shaded}

\begin{verbatim}
#> GAM — RMSE: 59275.36 | R²: 0.745
\end{verbatim}

The GAM provides the most interpretable nonlinear model, each smooth
term \(\hat{f}_j(x_j)\) can be plotted individually to show its
contribution to the prediction. The
\texttt{ti(median\_income,\ latitude)} tensor term reveals that the
income premium on house values is \textbf{steepest in the Bay Area} (lat
≈ 37--38°), a finding invisible to additive models.

\begin{center}\rule{0.5\linewidth}{0.5pt}\end{center}

\section{Model 5 --- Regression Splines (B-splines \& Natural
Splines)}\label{splines}

Regression splines fit piecewise polynomials with continuity constraints
at knots:

\[y_i = \beta_0 + \sum_{j=1}^{K+3} \beta_j b_j(x_i) + \varepsilon_i\]

A \textbf{cubic B-spline} with \(K\) knots has \(K + 4\) degrees of
freedom. A \textbf{natural spline} adds boundary linearity constraints,
reducing instability at the tails.

\begin{Shaded}
\begin{Highlighting}[]
\CommentTok{\# B{-}spline model }
\NormalTok{bspline\_model }\OtherTok{\textless{}{-}} \FunctionTok{lm}\NormalTok{(}
\NormalTok{  median\_house\_value }\SpecialCharTok{\textasciitilde{}}
    \FunctionTok{bs}\NormalTok{(median\_income,            }\AttributeTok{df =} \DecValTok{8}\NormalTok{) }\SpecialCharTok{+}   \CommentTok{\# 6 df ≈ 3 interior knots}
    \FunctionTok{bs}\NormalTok{(latitude,                 }\AttributeTok{df =} \DecValTok{8}\NormalTok{) }\SpecialCharTok{+}   \CommentTok{\# 8 df — bimodal, needs flexibility}
    \FunctionTok{bs}\NormalTok{(longitude,                }\AttributeTok{df =} \DecValTok{8}\NormalTok{) }\SpecialCharTok{+}
    \FunctionTok{bs}\NormalTok{(housing\_median\_age,       }\AttributeTok{df =} \DecValTok{8}\NormalTok{) }\SpecialCharTok{+}
    \FunctionTok{bs}\NormalTok{(rooms\_per\_household,      }\AttributeTok{df =} \DecValTok{7}\NormalTok{) }\SpecialCharTok{+}
    \FunctionTok{bs}\NormalTok{(population\_per\_household, }\AttributeTok{df =} \DecValTok{8}\NormalTok{) }\SpecialCharTok{+}
    \FunctionTok{bs}\NormalTok{(bedrooms\_per\_room,        }\AttributeTok{df =} \DecValTok{5}\NormalTok{),}
  \AttributeTok{data =}\NormalTok{ train\_sc}
\NormalTok{)}

\NormalTok{bspline\_pred }\OtherTok{\textless{}{-}} \FunctionTok{predict}\NormalTok{(bspline\_model, }\AttributeTok{newdata =}\NormalTok{ test\_sc)}

\CommentTok{\# Natural spline model }
\NormalTok{nspline\_model }\OtherTok{\textless{}{-}} \FunctionTok{lm}\NormalTok{(}
\NormalTok{  median\_house\_value }\SpecialCharTok{\textasciitilde{}}
    \FunctionTok{ns}\NormalTok{(median\_income,            }\AttributeTok{df =} \DecValTok{8}\NormalTok{) }\SpecialCharTok{+}
    \FunctionTok{ns}\NormalTok{(latitude,                 }\AttributeTok{df =} \DecValTok{8}\NormalTok{) }\SpecialCharTok{+}
    \FunctionTok{ns}\NormalTok{(longitude,                }\AttributeTok{df =} \DecValTok{8}\NormalTok{) }\SpecialCharTok{+}
    \FunctionTok{ns}\NormalTok{(housing\_median\_age,       }\AttributeTok{df =} \DecValTok{8}\NormalTok{) }\SpecialCharTok{+}
    \FunctionTok{ns}\NormalTok{(rooms\_per\_household,      }\AttributeTok{df =} \DecValTok{7}\NormalTok{) }\SpecialCharTok{+}
    \FunctionTok{ns}\NormalTok{(population\_per\_household, }\AttributeTok{df =} \DecValTok{8}\NormalTok{) }\SpecialCharTok{+}
    \FunctionTok{ns}\NormalTok{(bedrooms\_per\_room,        }\AttributeTok{df =} \DecValTok{5}\NormalTok{),}
  \AttributeTok{data =}\NormalTok{ train\_sc}
\NormalTok{)}

\NormalTok{nspline\_pred }\OtherTok{\textless{}{-}} \FunctionTok{predict}\NormalTok{(nspline\_model, }\AttributeTok{newdata =}\NormalTok{ test\_sc)}

\FunctionTok{cat}\NormalTok{(}\StringTok{"B{-}Spline       — RMSE:"}\NormalTok{, }\FunctionTok{round}\NormalTok{(}\FunctionTok{rmse}\NormalTok{(y\_test, bspline\_pred), }\DecValTok{2}\NormalTok{),}
    \StringTok{"| R²:"}\NormalTok{, }\FunctionTok{round}\NormalTok{(}\FunctionTok{cor}\NormalTok{(y\_test, bspline\_pred)}\SpecialCharTok{\^{}}\DecValTok{2}\NormalTok{, }\DecValTok{4}\NormalTok{), }\StringTok{"}\SpecialCharTok{\textbackslash{}n}\StringTok{"}\NormalTok{)}
\end{Highlighting}
\end{Shaded}

\begin{verbatim}
#> B-Spline       — RMSE: 61987.56 | R²: 0.7211
\end{verbatim}

\begin{Shaded}
\begin{Highlighting}[]
\FunctionTok{cat}\NormalTok{(}\StringTok{"Natural Spline — RMSE:"}\NormalTok{, }\FunctionTok{round}\NormalTok{(}\FunctionTok{rmse}\NormalTok{(y\_test, nspline\_pred), }\DecValTok{2}\NormalTok{),}
    \StringTok{"| R²:"}\NormalTok{, }\FunctionTok{round}\NormalTok{(}\FunctionTok{cor}\NormalTok{(y\_test, nspline\_pred)}\SpecialCharTok{\^{}}\DecValTok{2}\NormalTok{, }\DecValTok{4}\NormalTok{), }\StringTok{"}\SpecialCharTok{\textbackslash{}n}\StringTok{"}\NormalTok{)}
\end{Highlighting}
\end{Shaded}

\begin{verbatim}
#> Natural Spline — RMSE: 61365.27 | R²: 0.7267
\end{verbatim}

\begin{Shaded}
\begin{Highlighting}[]
\CommentTok{\# Compare linear, B{-}spline, natural spline fits for median\_income}
\NormalTok{grid\_df }\OtherTok{\textless{}{-}} \FunctionTok{data.frame}\NormalTok{(}
  \AttributeTok{median\_income            =} \FunctionTok{seq}\NormalTok{(}\FunctionTok{min}\NormalTok{(train\_sc}\SpecialCharTok{$}\NormalTok{median\_income), }\FunctionTok{max}\NormalTok{(train\_sc}\SpecialCharTok{$}\NormalTok{median\_income), }\AttributeTok{length =} \DecValTok{300}\NormalTok{),}
  \AttributeTok{latitude                 =} \FunctionTok{median}\NormalTok{(train\_sc}\SpecialCharTok{$}\NormalTok{latitude),}
  \AttributeTok{longitude                =} \FunctionTok{median}\NormalTok{(train\_sc}\SpecialCharTok{$}\NormalTok{longitude),}
  \AttributeTok{housing\_median\_age       =} \FunctionTok{median}\NormalTok{(train\_sc}\SpecialCharTok{$}\NormalTok{housing\_median\_age),}
  \AttributeTok{rooms\_per\_household      =} \FunctionTok{median}\NormalTok{(train\_sc}\SpecialCharTok{$}\NormalTok{rooms\_per\_household),}
  \AttributeTok{bedrooms\_per\_room        =} \FunctionTok{median}\NormalTok{(train\_sc}\SpecialCharTok{$}\NormalTok{bedrooms\_per\_room),}
  \AttributeTok{population\_per\_household =} \FunctionTok{median}\NormalTok{(train\_sc}\SpecialCharTok{$}\NormalTok{population\_per\_household)}
  
\NormalTok{)}

\NormalTok{grid\_df}\SpecialCharTok{$}\NormalTok{lm\_fit       }\OtherTok{\textless{}{-}} \FunctionTok{predict}\NormalTok{(lm\_model,      }\AttributeTok{newdata =}\NormalTok{ grid\_df }\SpecialCharTok{\%\textgreater{}\%} \FunctionTok{mutate}\NormalTok{(}\AttributeTok{households =} \FunctionTok{median}\NormalTok{(train\_sc}\SpecialCharTok{$}\NormalTok{households)))}
\NormalTok{grid\_df}\SpecialCharTok{$}\NormalTok{bs\_fit       }\OtherTok{\textless{}{-}} \FunctionTok{predict}\NormalTok{(bspline\_model,  }\AttributeTok{newdata =}\NormalTok{ grid\_df)}
\NormalTok{grid\_df}\SpecialCharTok{$}\NormalTok{ns\_fit       }\OtherTok{\textless{}{-}} \FunctionTok{predict}\NormalTok{(nspline\_model,  }\AttributeTok{newdata =}\NormalTok{ grid\_df)}

\NormalTok{grid\_long }\OtherTok{\textless{}{-}}\NormalTok{ grid\_df }\SpecialCharTok{\%\textgreater{}\%}
  \FunctionTok{pivot\_longer}\NormalTok{(}\FunctionTok{c}\NormalTok{(lm\_fit, bs\_fit, ns\_fit), }\AttributeTok{names\_to =} \StringTok{"model"}\NormalTok{, }\AttributeTok{values\_to =} \StringTok{"fitted"}\NormalTok{) }\SpecialCharTok{\%\textgreater{}\%}
  \FunctionTok{mutate}\NormalTok{(}\AttributeTok{model =} \FunctionTok{recode}\NormalTok{(model,}
    \AttributeTok{lm\_fit =} \StringTok{"Linear"}\NormalTok{, }\AttributeTok{bs\_fit =} \StringTok{"B{-}Spline (df=6)"}\NormalTok{, }\AttributeTok{ns\_fit =} \StringTok{"Natural Spline (df=6)"}\NormalTok{))}

\FunctionTok{ggplot}\NormalTok{(grid\_long, }\FunctionTok{aes}\NormalTok{(}\AttributeTok{x =}\NormalTok{ median\_income, }\AttributeTok{y =}\NormalTok{ fitted, }\AttributeTok{color =}\NormalTok{ model, }\AttributeTok{linetype =}\NormalTok{ model)) }\SpecialCharTok{+}
  \FunctionTok{geom\_line}\NormalTok{(}\AttributeTok{linewidth =} \FloatTok{1.1}\NormalTok{) }\SpecialCharTok{+}
  \FunctionTok{scale\_color\_manual}\NormalTok{(}\AttributeTok{values =} \FunctionTok{c}\NormalTok{(}\StringTok{"Linear"} \OtherTok{=} \StringTok{"\#7f8c8d"}\NormalTok{,}
                                \StringTok{"B{-}Spline (df=6)"} \OtherTok{=} \StringTok{"\#2e6da4"}\NormalTok{,}
                                \StringTok{"Natural Spline (df=6)"} \OtherTok{=} \StringTok{"\#c0392b"}\NormalTok{)) }\SpecialCharTok{+}
  \FunctionTok{labs}\NormalTok{(}
    \AttributeTok{title    =} \StringTok{"Regression Splines: Fitted Curves for Median Income"}\NormalTok{,}
    \AttributeTok{subtitle =} \StringTok{"Splines capture nonlinearity; natural spline is more stable at boundaries"}\NormalTok{,}
    \AttributeTok{x =} \StringTok{"Median Income (scaled)"}\NormalTok{, }\AttributeTok{y =} \StringTok{"Fitted House Value ($)"}\NormalTok{,}
    \AttributeTok{color =} \StringTok{"Model"}\NormalTok{, }\AttributeTok{linetype =} \StringTok{"Model"}
\NormalTok{  )}
\end{Highlighting}
\end{Shaded}

\pandocbounded{\includegraphics[keepaspectratio]{California_Housing_Analysis_files/figure-latex/spline-plot-1.pdf}}

\begin{Shaded}
\begin{Highlighting}[]
\NormalTok{all\_preds[[}\StringTok{"BSpline"}\NormalTok{]]       }\OtherTok{\textless{}{-}}\NormalTok{ bspline\_pred}
\NormalTok{all\_preds[[}\StringTok{"NaturalSpline"}\NormalTok{]] }\OtherTok{\textless{}{-}}\NormalTok{ nspline\_pred}

\NormalTok{all\_metrics[[}\StringTok{"BSpline"}\NormalTok{]]       }\OtherTok{\textless{}{-}} \FunctionTok{calc\_metrics}\NormalTok{(y\_test, bspline\_pred, }\StringTok{"4a. B{-}Spline Regression"}\NormalTok{)}
\NormalTok{all\_metrics[[}\StringTok{"NaturalSpline"}\NormalTok{]] }\OtherTok{\textless{}{-}} \FunctionTok{calc\_metrics}\NormalTok{(y\_test, nspline\_pred, }\StringTok{"4b. Natural Spline Regression"}\NormalTok{)}
\end{Highlighting}
\end{Shaded}

Both spline models substantially outperform the linear baseline and the
step function. The natural spline is preferred in practice because it
imposes linearity at the boundaries, avoiding erratic tail behaviour.

\begin{center}\rule{0.5\linewidth}{0.5pt}\end{center}

\section{\texorpdfstring{Model 6 --- Smoothing Spline (via GAM
\texttt{s()})}{Model 6 --- Smoothing Spline (via GAM s())}}\label{smoothing-spline}

Smoothing splines minimise the penalised RSS:

\[\sum_{i=1}^{n}(y_i - g(x_i))^2 + \lambda \int g''(t)^2 \, dt\]

The tuning parameter \(\lambda\) controls the bias-variance trade-off.
Large \(\lambda\) → smooth \(g\); \(\lambda \to 0\) → interpolates the
data. \(\lambda\) is selected automatically via \textbf{REML}
cross-validation when using \texttt{mgcv::gam()}.

\begin{Shaded}
\begin{Highlighting}[]
\CommentTok{\# Single{-}predictor smoothing spline for median\_income (illustrative)}
\NormalTok{ss\_model\_1d }\OtherTok{\textless{}{-}}\NormalTok{ mgcv}\SpecialCharTok{::}\FunctionTok{gam}\NormalTok{(}
\NormalTok{  median\_house\_value }\SpecialCharTok{\textasciitilde{}} \FunctionTok{s}\NormalTok{(median\_income, }\AttributeTok{k =} \DecValTok{10}\NormalTok{, }\AttributeTok{bs =} \StringTok{"cr"}\NormalTok{),  }\CommentTok{\# cubic regression spline}
  \AttributeTok{data   =}\NormalTok{ train\_sc,}
  \AttributeTok{method =} \StringTok{"REML"}
\NormalTok{)}

\FunctionTok{summary}\NormalTok{(ss\_model\_1d)}
\end{Highlighting}
\end{Shaded}

\begin{verbatim}
#> 
#> Family: gaussian 
#> Link function: identity 
#> 
#> Formula:
#> median_house_value ~ s(median_income, k = 10, bs = "cr")
#> 
#> Parametric coefficients:
#>             Estimate Std. Error t value Pr(>|t|)    
#> (Intercept) 206427.6      639.3   322.9   <2e-16 ***
#> ---
#> Signif. codes:  0 '***' 0.001 '**' 0.01 '*' 0.05 '.' 0.1 ' ' 1
#> 
#> Approximate significance of smooth terms:
#>                    edf Ref.df    F p-value    
#> s(median_income) 8.112  8.726 1809  <2e-16 ***
#> ---
#> Signif. codes:  0 '***' 0.001 '**' 0.01 '*' 0.05 '.' 0.1 ' ' 1
#> 
#> R-sq.(adj) =  0.489   Deviance explained = 48.9%
#> -REML = 2.1029e+05  Scale est. = 6.7482e+09  n = 16513
\end{verbatim}

\begin{Shaded}
\begin{Highlighting}[]
\FunctionTok{cat}\NormalTok{(}\StringTok{"}\SpecialCharTok{\textbackslash{}n}\StringTok{Effective df (median\_income smooth):"}\NormalTok{,}
    \FunctionTok{round}\NormalTok{(}\FunctionTok{sum}\NormalTok{(ss\_model\_1d}\SpecialCharTok{$}\NormalTok{edf), }\DecValTok{2}\NormalTok{), }\StringTok{"}\SpecialCharTok{\textbackslash{}n}\StringTok{"}\NormalTok{)}
\end{Highlighting}
\end{Shaded}

\begin{verbatim}
#> 
#> Effective df (median_income smooth): 9.11
\end{verbatim}

\begin{Shaded}
\begin{Highlighting}[]
\CommentTok{\# Plot the smoothing spline with uncertainty band}
\FunctionTok{plot}\NormalTok{(ss\_model\_1d, }\AttributeTok{shade =} \ConstantTok{TRUE}\NormalTok{, }\AttributeTok{col =} \StringTok{"\#2e6da4"}\NormalTok{, }\AttributeTok{shade.col =} \StringTok{"lightblue"}\NormalTok{,}
     \AttributeTok{seWithMean =} \ConstantTok{TRUE}\NormalTok{, }\AttributeTok{residuals =} \ConstantTok{TRUE}\NormalTok{, }\AttributeTok{pch =} \DecValTok{20}\NormalTok{, }\AttributeTok{cex =} \FloatTok{0.1}\NormalTok{,}
     \AttributeTok{main =} \StringTok{"Smoothing Spline: s(median\_income) — REML{-}selected λ"}\NormalTok{,}
     \AttributeTok{xlab =} \StringTok{"Median Income (scaled)"}\NormalTok{, }\AttributeTok{ylab =} \StringTok{"Effect on House Value ($)"}\NormalTok{)}
\FunctionTok{abline}\NormalTok{(}\AttributeTok{h =} \DecValTok{0}\NormalTok{, }\AttributeTok{lty =} \DecValTok{2}\NormalTok{, }\AttributeTok{col =} \StringTok{"grey50"}\NormalTok{)}
\end{Highlighting}
\end{Shaded}

\pandocbounded{\includegraphics[keepaspectratio]{California_Housing_Analysis_files/figure-latex/smoothing-spline-plot-1.pdf}}

\begin{Shaded}
\begin{Highlighting}[]
\CommentTok{\# Multi{-}predictor smoothing spline model}
\NormalTok{ss\_multi }\OtherTok{\textless{}{-}}\NormalTok{ mgcv}\SpecialCharTok{::}\FunctionTok{gam}\NormalTok{(}
\NormalTok{  median\_house\_value }\SpecialCharTok{\textasciitilde{}}
    \FunctionTok{s}\NormalTok{(median\_income,            }\AttributeTok{k =} \DecValTok{10}\NormalTok{, }\AttributeTok{bs =} \StringTok{"cr"}\NormalTok{) }\SpecialCharTok{+}
    \FunctionTok{s}\NormalTok{(latitude,                 }\AttributeTok{k =} \DecValTok{12}\NormalTok{, }\AttributeTok{bs =} \StringTok{"cr"}\NormalTok{) }\SpecialCharTok{+}
    \FunctionTok{s}\NormalTok{(longitude,                }\AttributeTok{k =} \DecValTok{12}\NormalTok{, }\AttributeTok{bs =} \StringTok{"cr"}\NormalTok{) }\SpecialCharTok{+}
    \FunctionTok{s}\NormalTok{(housing\_median\_age,       }\AttributeTok{k =}  \DecValTok{8}\NormalTok{, }\AttributeTok{bs =} \StringTok{"cr"}\NormalTok{) }\SpecialCharTok{+}
    \FunctionTok{s}\NormalTok{(rooms\_per\_household,      }\AttributeTok{k =}  \DecValTok{8}\NormalTok{, }\AttributeTok{bs =} \StringTok{"cr"}\NormalTok{) }\SpecialCharTok{+}
    \FunctionTok{s}\NormalTok{(population\_per\_household, }\AttributeTok{k =}  \DecValTok{8}\NormalTok{, }\AttributeTok{bs =} \StringTok{"cr"}\NormalTok{) }\SpecialCharTok{+}
\NormalTok{    bedrooms\_per\_room,}
  \AttributeTok{data   =}\NormalTok{ train\_sc,}
  \AttributeTok{method =} \StringTok{"REML"}
\NormalTok{)}

\NormalTok{ss\_pred }\OtherTok{\textless{}{-}} \FunctionTok{predict}\NormalTok{(ss\_multi, }\AttributeTok{newdata =}\NormalTok{ test\_sc)}

\NormalTok{all\_preds[[}\StringTok{"SmoothingSpline"}\NormalTok{]] }\OtherTok{\textless{}{-}}\NormalTok{ ss\_pred}
\NormalTok{all\_metrics[[}\StringTok{"SmoothingSpline"}\NormalTok{]] }\OtherTok{\textless{}{-}} \FunctionTok{calc\_metrics}\NormalTok{(y\_test, ss\_pred, }\StringTok{"5. Smoothing Spline (REML)"}\NormalTok{)}

\FunctionTok{cat}\NormalTok{(}\StringTok{"Smoothing Spline — RMSE:"}\NormalTok{, all\_metrics}\SpecialCharTok{$}\NormalTok{SmoothingSpline}\SpecialCharTok{$}\NormalTok{RMSE,}
    \StringTok{"| R²:"}\NormalTok{, all\_metrics}\SpecialCharTok{$}\NormalTok{SmoothingSpline}\SpecialCharTok{$}\NormalTok{R2, }\StringTok{"}\SpecialCharTok{\textbackslash{}n}\StringTok{"}\NormalTok{)}
\end{Highlighting}
\end{Shaded}

\begin{verbatim}
#> Smoothing Spline — RMSE: 61351.49 | R²: 0.7268
\end{verbatim}

\begin{Shaded}
\begin{Highlighting}[]
\FunctionTok{par}\NormalTok{(}\AttributeTok{mfrow =} \FunctionTok{c}\NormalTok{(}\DecValTok{3}\NormalTok{, }\DecValTok{2}\NormalTok{))}
\FunctionTok{plot}\NormalTok{(ss\_multi, }\AttributeTok{pages =} \DecValTok{0}\NormalTok{, }\AttributeTok{shade =} \ConstantTok{TRUE}\NormalTok{, }\AttributeTok{col =} \StringTok{"\#c0392b"}\NormalTok{,}
     \AttributeTok{shade.col =} \StringTok{"\#f5c6c6"}\NormalTok{, }\AttributeTok{seWithMean =} \ConstantTok{TRUE}\NormalTok{,}
     \AttributeTok{main =} \StringTok{"Smoothing Spline Terms"}\NormalTok{)}
\end{Highlighting}
\end{Shaded}

\pandocbounded{\includegraphics[keepaspectratio]{California_Housing_Analysis_files/figure-latex/smoothing-spline-plots-multi-1.pdf}}

\begin{Shaded}
\begin{Highlighting}[]
\FunctionTok{par}\NormalTok{(}\AttributeTok{mfrow =} \FunctionTok{c}\NormalTok{(}\DecValTok{1}\NormalTok{, }\DecValTok{1}\NormalTok{))}
\end{Highlighting}
\end{Shaded}

The smoothing spline for \texttt{median\_income} reveals a clear
S-shaped relationship, steepest around income levels 2--5 (scaled),
flattening for very high incomes. The \texttt{latitude} smooth shows the
bimodal pattern. The effective degrees of freedom for each smooth are
selected automatically, a major advantage over manually choosing knot
placements.

\begin{center}\rule{0.5\linewidth}{0.5pt}\end{center}

\section{Model 7 --- Local Regression (LOESS)}\label{local-reg}

Local regression fits weighted least-squares models in a moving
neighbourhood of each point. The span parameter \(\alpha\) controls the
neighbourhood width, analogous to \(\lambda\) in smoothing splines.

\[\hat{f}(x_0) = \arg\min_{\beta} \sum_{i=1}^{n} K\!\left(\frac{x_i - x_0}{h}\right)(y_i - \beta_0 - \beta_1 x_i)^2\]

For multiple predictors we use \textbf{np} (nonparametric) kernel
regression with bandwidth selected by cross-validation. For speed, we
fit LOESS on a sample and use a GAM \texttt{lo()} approximation for the
full model.

\begin{Shaded}
\begin{Highlighting}[]
\CommentTok{\# ── Illustrative 1D LOESS: income vs house value ──────────────────────────────}
\NormalTok{loess\_1d }\OtherTok{\textless{}{-}} \FunctionTok{loess}\NormalTok{(}
\NormalTok{  median\_house\_value }\SpecialCharTok{\textasciitilde{}}\NormalTok{ median\_income,}
  \AttributeTok{data =}\NormalTok{ train\_sc, }\AttributeTok{span =} \FloatTok{0.25}\NormalTok{, }\AttributeTok{degree =} \DecValTok{2}
\NormalTok{)}

\NormalTok{loess\_grid }\OtherTok{\textless{}{-}} \FunctionTok{data.frame}\NormalTok{(}\AttributeTok{median\_income =} \FunctionTok{seq}\NormalTok{(}
  \FunctionTok{min}\NormalTok{(train\_sc}\SpecialCharTok{$}\NormalTok{median\_income), }\FunctionTok{max}\NormalTok{(train\_sc}\SpecialCharTok{$}\NormalTok{median\_income), }\AttributeTok{length =} \DecValTok{200}\NormalTok{))}
\NormalTok{loess\_grid}\SpecialCharTok{$}\NormalTok{fit }\OtherTok{\textless{}{-}} \FunctionTok{predict}\NormalTok{(loess\_1d, }\AttributeTok{newdata =}\NormalTok{ loess\_grid)}

\FunctionTok{ggplot}\NormalTok{(train\_sc, }\FunctionTok{aes}\NormalTok{(}\AttributeTok{x =}\NormalTok{ median\_income, }\AttributeTok{y =}\NormalTok{ median\_house\_value)) }\SpecialCharTok{+}
  \FunctionTok{geom\_point}\NormalTok{(}\AttributeTok{alpha =} \FloatTok{0.06}\NormalTok{, }\AttributeTok{color =} \StringTok{"\#2e6da4"}\NormalTok{, }\AttributeTok{size =} \FloatTok{0.5}\NormalTok{) }\SpecialCharTok{+}
  \FunctionTok{geom\_line}\NormalTok{(}\AttributeTok{data =}\NormalTok{ loess\_grid, }\FunctionTok{aes}\NormalTok{(}\AttributeTok{y =}\NormalTok{ fit), }\AttributeTok{color =} \StringTok{"\#c0392b"}\NormalTok{, }\AttributeTok{linewidth =} \FloatTok{1.2}\NormalTok{) }\SpecialCharTok{+}
  \FunctionTok{labs}\NormalTok{(}
    \AttributeTok{title    =} \StringTok{"Local Regression (LOESS, span = 0.25): Median Income vs House Value"}\NormalTok{,}
    \AttributeTok{subtitle =} \StringTok{"Locally{-}weighted polynomial fit — no global functional form assumed"}\NormalTok{,}
    \AttributeTok{x =} \StringTok{"Median Income (scaled)"}\NormalTok{, }\AttributeTok{y =} \StringTok{"House Value ($)"}
\NormalTok{  )}
\end{Highlighting}
\end{Shaded}

\pandocbounded{\includegraphics[keepaspectratio]{California_Housing_Analysis_files/figure-latex/loess-illustrative-1.pdf}}

\begin{Shaded}
\begin{Highlighting}[]
\CommentTok{\# ── Multivariate local regression via GAM lo() ────────────────────────────────}
\CommentTok{\# mgcv::gam with lo() terms provides a practical multivariate LOESS approximation}
\NormalTok{loess\_model }\OtherTok{\textless{}{-}}\NormalTok{mgcv}\SpecialCharTok{::}\FunctionTok{gam}\NormalTok{(}
\NormalTok{  median\_house\_value }\SpecialCharTok{\textasciitilde{}}
\NormalTok{    gam}\SpecialCharTok{::}\FunctionTok{lo}\NormalTok{(median\_income, }\AttributeTok{span =} \FloatTok{0.3}\NormalTok{) }\SpecialCharTok{+}
\NormalTok{    gam}\SpecialCharTok{::}\FunctionTok{lo}\NormalTok{(latitude,      }\AttributeTok{span =} \FloatTok{0.3}\NormalTok{) }\SpecialCharTok{+}
\NormalTok{    gam}\SpecialCharTok{::}\FunctionTok{lo}\NormalTok{(longitude,     }\AttributeTok{span =} \FloatTok{0.3}\NormalTok{) }\SpecialCharTok{+}
\NormalTok{    gam}\SpecialCharTok{::}\FunctionTok{lo}\NormalTok{(housing\_median\_age, }\AttributeTok{span =} \FloatTok{0.4}\NormalTok{) }\SpecialCharTok{+}
\NormalTok{    gam}\SpecialCharTok{::}\FunctionTok{lo}\NormalTok{(rooms\_per\_household, }\AttributeTok{span =} \FloatTok{0.4}\NormalTok{) }\SpecialCharTok{+}
\NormalTok{    rooms\_per\_household }\SpecialCharTok{+}\NormalTok{ bedrooms\_per\_room }\SpecialCharTok{+}
\NormalTok{    population\_per\_household,}
  \AttributeTok{data =}\NormalTok{ train\_sc}
\NormalTok{)}

\NormalTok{loess\_pred }\OtherTok{\textless{}{-}} \FunctionTok{predict}\NormalTok{(loess\_model, }\AttributeTok{newdata =}\NormalTok{ test\_sc)}

\NormalTok{all\_preds[[}\StringTok{"LocalRegression"}\NormalTok{]] }\OtherTok{\textless{}{-}}\NormalTok{ loess\_pred}
\NormalTok{all\_metrics[[}\StringTok{"LocalRegression"}\NormalTok{]] }\OtherTok{\textless{}{-}} \FunctionTok{calc\_metrics}\NormalTok{(y\_test, loess\_pred, }\StringTok{"6. Local Regression (LOESS)"}\NormalTok{)}

\FunctionTok{cat}\NormalTok{(}\StringTok{"Local Regression — RMSE:"}\NormalTok{, all\_metrics}\SpecialCharTok{$}\NormalTok{LocalRegression}\SpecialCharTok{$}\NormalTok{RMSE,}
    \StringTok{"| R²:"}\NormalTok{, all\_metrics}\SpecialCharTok{$}\NormalTok{LocalRegression}\SpecialCharTok{$}\NormalTok{R2, }\StringTok{"}\SpecialCharTok{\textbackslash{}n}\StringTok{"}\NormalTok{)}
\end{Highlighting}
\end{Shaded}

\begin{verbatim}
#> Local Regression — RMSE: 73005.27 | R²: 0.6134
\end{verbatim}

LOESS is fully nonparametric, it makes no assumption about the global
functional form. However, it does not extrapolate well and is
computationally expensive. For our 11-predictor problem, GAM provides
more scalable alternative.

\section{Model 8 --- Random Forest}\label{rf}

Random Forest is an ensemble of decorrelated decision trees trained on
bootstrap samples. Each tree is grown on a random subset of features
(\texttt{mtry}) and predictions are averaged. It implicitly captures all
nonlinearities and interactions without explicit specification ---
analogous to a fully nonparametric model.

\begin{Shaded}
\begin{Highlighting}[]
\FunctionTok{set.seed}\NormalTok{(}\DecValTok{42}\NormalTok{)}
\NormalTok{rf\_model }\OtherTok{\textless{}{-}} \FunctionTok{randomForest}\NormalTok{(}
  \AttributeTok{x         =}\NormalTok{ X\_train,}
  \AttributeTok{y         =}\NormalTok{ y\_train,}
  \AttributeTok{ntree     =} \DecValTok{100}\NormalTok{,}
  \AttributeTok{mtry      =} \DecValTok{4}\NormalTok{,          }\CommentTok{\# ≈ sqrt(p) for regression}
  \AttributeTok{importance =} \ConstantTok{TRUE}\NormalTok{,}
  \AttributeTok{do.trace   =} \ConstantTok{FALSE}
\NormalTok{)}

\FunctionTok{print}\NormalTok{(rf\_model)}
\end{Highlighting}
\end{Shaded}

\begin{verbatim}
#> 
#> Call:
#>  randomForest(x = X_train, y = y_train, ntree = 100, mtry = 4,      importance = TRUE, do.trace = FALSE) 
#>                Type of random forest: regression
#>                      Number of trees: 100
#> No. of variables tried at each split: 4
#> 
#>           Mean of squared residuals: 2516742507
#>                     % Var explained: 80.93
\end{verbatim}

\begin{Shaded}
\begin{Highlighting}[]
\CommentTok{\# Variable importance}
\NormalTok{imp\_df }\OtherTok{\textless{}{-}} \FunctionTok{as.data.frame}\NormalTok{(}\FunctionTok{importance}\NormalTok{(rf\_model)) }\SpecialCharTok{\%\textgreater{}\%}
  \FunctionTok{rownames\_to\_column}\NormalTok{(}\StringTok{"Feature"}\NormalTok{) }\SpecialCharTok{\%\textgreater{}\%}
  \FunctionTok{arrange}\NormalTok{(}\FunctionTok{desc}\NormalTok{(}\StringTok{\textasciigrave{}}\AttributeTok{\%IncMSE}\StringTok{\textasciigrave{}}\NormalTok{))}

\FunctionTok{ggplot}\NormalTok{(imp\_df, }\FunctionTok{aes}\NormalTok{(}\AttributeTok{x =} \FunctionTok{reorder}\NormalTok{(Feature, }\StringTok{\textasciigrave{}}\AttributeTok{\%IncMSE}\StringTok{\textasciigrave{}}\NormalTok{), }\AttributeTok{y =} \StringTok{\textasciigrave{}}\AttributeTok{\%IncMSE}\StringTok{\textasciigrave{}}\NormalTok{)) }\SpecialCharTok{+}
  \FunctionTok{geom\_col}\NormalTok{(}\AttributeTok{fill =} \StringTok{"\#2e6da4"}\NormalTok{, }\AttributeTok{alpha =} \FloatTok{0.85}\NormalTok{) }\SpecialCharTok{+}
  \FunctionTok{coord\_flip}\NormalTok{() }\SpecialCharTok{+}
  \FunctionTok{labs}\NormalTok{(}
    \AttributeTok{title =} \StringTok{"Random Forest — Variable Importance (\% Increase in MSE)"}\NormalTok{,}
    \AttributeTok{x =} \ConstantTok{NULL}\NormalTok{, }\AttributeTok{y =} \StringTok{"\% Increase in MSE if feature permuted"}
\NormalTok{  )}
\end{Highlighting}
\end{Shaded}

\pandocbounded{\includegraphics[keepaspectratio]{California_Housing_Analysis_files/figure-latex/rf-importance-plot-1.pdf}}

\begin{Shaded}
\begin{Highlighting}[]
\CommentTok{\# Partial Dependence Plots for top 2 features}
\NormalTok{pd\_income }\OtherTok{\textless{}{-}} \FunctionTok{partial}\NormalTok{(rf\_model, }\AttributeTok{pred.var =} \StringTok{"median\_income"}\NormalTok{,}
                     \AttributeTok{train =}\NormalTok{ X\_train, }\AttributeTok{grid.resolution =} \DecValTok{30}\NormalTok{)}
\NormalTok{pd\_population\_per\_household    }\OtherTok{\textless{}{-}} \FunctionTok{partial}\NormalTok{(rf\_model, }\AttributeTok{pred.var =} \StringTok{"population\_per\_household"}\NormalTok{,}
                     \AttributeTok{train =}\NormalTok{ X\_train, }\AttributeTok{grid.resolution =} \DecValTok{30}\NormalTok{)}

\NormalTok{pd\_lat }\OtherTok{\textless{}{-}} \FunctionTok{partial}\NormalTok{(rf\_model, }\AttributeTok{pred.var =} \StringTok{"latitude"}\NormalTok{,}
                     \AttributeTok{train =}\NormalTok{ X\_train, }\AttributeTok{grid.resolution =} \DecValTok{30}\NormalTok{)}

\NormalTok{p\_income }\OtherTok{\textless{}{-}} \FunctionTok{autoplot}\NormalTok{(pd\_income, }\AttributeTok{smooth =} \ConstantTok{TRUE}\NormalTok{) }\SpecialCharTok{+}
  \FunctionTok{labs}\NormalTok{(}\AttributeTok{title =} \StringTok{"PDP: Median Income (RF)"}\NormalTok{, }\AttributeTok{x =} \StringTok{"Income (scaled)"}\NormalTok{, }\AttributeTok{y =} \StringTok{"Fitted Value"}\NormalTok{) }\SpecialCharTok{+}
  \FunctionTok{theme\_minimal}\NormalTok{()}

\NormalTok{p\_lat }\OtherTok{\textless{}{-}} \FunctionTok{autoplot}\NormalTok{(pd\_lat, }\AttributeTok{smooth =} \ConstantTok{TRUE}\NormalTok{) }\SpecialCharTok{+}
  \FunctionTok{labs}\NormalTok{(}\AttributeTok{title =} \StringTok{"PDP: Population\_per\_household (RF)"}\NormalTok{, }\AttributeTok{x =} \StringTok{"Population\_per\_household (scaled)"}\NormalTok{, }\AttributeTok{y =} \StringTok{"Fitted Value"}\NormalTok{) }\SpecialCharTok{+}
  \FunctionTok{theme\_minimal}\NormalTok{()}

\NormalTok{p\_income }\SpecialCharTok{+}\NormalTok{ p\_lat}
\end{Highlighting}
\end{Shaded}

\pandocbounded{\includegraphics[keepaspectratio]{California_Housing_Analysis_files/figure-latex/rf-pdp-1.pdf}}

\begin{Shaded}
\begin{Highlighting}[]
\NormalTok{rf\_pred }\OtherTok{\textless{}{-}} \FunctionTok{predict}\NormalTok{(rf\_model, }\AttributeTok{newdata =}\NormalTok{ X\_test)}

\NormalTok{all\_preds[[}\StringTok{"RandomForest"}\NormalTok{]] }\OtherTok{\textless{}{-}}\NormalTok{ rf\_pred}
\NormalTok{all\_metrics[[}\StringTok{"RandomForest"}\NormalTok{]] }\OtherTok{\textless{}{-}} \FunctionTok{calc\_metrics}\NormalTok{(y\_test, rf\_pred, }\StringTok{"8. Random Forest"}\NormalTok{)}

\FunctionTok{cat}\NormalTok{(}\StringTok{"Random Forest — RMSE:"}\NormalTok{, all\_metrics}\SpecialCharTok{$}\NormalTok{RandomForest}\SpecialCharTok{$}\NormalTok{RMSE,}
    \StringTok{"| R²:"}\NormalTok{, all\_metrics}\SpecialCharTok{$}\NormalTok{RandomForest}\SpecialCharTok{$}\NormalTok{R2, }\StringTok{"}\SpecialCharTok{\textbackslash{}n}\StringTok{"}\NormalTok{)}
\end{Highlighting}
\end{Shaded}

\begin{verbatim}
#> Random Forest — RMSE: 50251.31 | R²: 0.8181
\end{verbatim}

Random Forest ranks \texttt{median\_income} first (\textasciitilde42\%
importance), followed by \texttt{population\ per\ household} and
\texttt{House\ median\ age}.

\begin{center}\rule{0.5\linewidth}{0.5pt}\end{center}

\section{Model 9 --- XGBoost}\label{xgboost}

XGBoost (eXtreme Gradient Boosting) sequentially fits shallow trees to
residuals using second-order gradient information. L1/L2 regularisation
and subsampling control overfitting --- playing the role of \(\lambda\)
in the smoothing spline penalty framework.

\begin{Shaded}
\begin{Highlighting}[]
\NormalTok{dtrain }\OtherTok{\textless{}{-}} \FunctionTok{xgb.DMatrix}\NormalTok{(}\AttributeTok{data =} \FunctionTok{as.matrix}\NormalTok{(X\_train), }\AttributeTok{label =}\NormalTok{ y\_train)}
\NormalTok{dtest  }\OtherTok{\textless{}{-}} \FunctionTok{xgb.DMatrix}\NormalTok{(}\AttributeTok{data =} \FunctionTok{as.matrix}\NormalTok{(X\_test),  }\AttributeTok{label =}\NormalTok{ y\_test)}

\NormalTok{xgb\_params }\OtherTok{\textless{}{-}} \FunctionTok{list}\NormalTok{(}
  \AttributeTok{objective        =} \StringTok{"reg:squarederror"}\NormalTok{,}
  \AttributeTok{eta              =} \FloatTok{0.05}\NormalTok{,}
  \AttributeTok{max\_depth        =} \DecValTok{6}\NormalTok{,}
  \AttributeTok{subsample        =} \FloatTok{0.8}\NormalTok{,}
  \AttributeTok{colsample\_bytree =} \FloatTok{0.8}\NormalTok{,}
  \AttributeTok{min\_child\_weight =} \DecValTok{5}\NormalTok{,}
  \AttributeTok{gamma            =} \FloatTok{0.1}
\NormalTok{)}

\FunctionTok{set.seed}\NormalTok{(}\DecValTok{42}\NormalTok{)}
\NormalTok{xgb\_model }\OtherTok{\textless{}{-}} \FunctionTok{xgb.train}\NormalTok{(}
  \AttributeTok{params                =}\NormalTok{ xgb\_params,}
  \AttributeTok{data                  =}\NormalTok{ dtrain,}
  \AttributeTok{nrounds               =} \DecValTok{500}\NormalTok{,}
  \AttributeTok{watchlist             =} \FunctionTok{list}\NormalTok{(}\AttributeTok{train =}\NormalTok{ dtrain, }\AttributeTok{test =}\NormalTok{ dtest),}
  \AttributeTok{early\_stopping\_rounds =} \DecValTok{30}\NormalTok{,}
  \AttributeTok{verbose               =} \DecValTok{0}
\NormalTok{)}

\FunctionTok{cat}\NormalTok{(}\StringTok{"Best iteration:"}\NormalTok{, xgb\_model}\SpecialCharTok{$}\NormalTok{best\_iteration, }\StringTok{"}\SpecialCharTok{\textbackslash{}n}\StringTok{"}\NormalTok{)}
\end{Highlighting}
\end{Shaded}

\begin{verbatim}
#> Best iteration:
\end{verbatim}

\begin{Shaded}
\begin{Highlighting}[]
\FunctionTok{cat}\NormalTok{(}\StringTok{"Best test RMSE:"}\NormalTok{, xgb\_model}\SpecialCharTok{$}\NormalTok{best\_score, }\StringTok{"}\SpecialCharTok{\textbackslash{}n}\StringTok{"}\NormalTok{)}
\end{Highlighting}
\end{Shaded}

\begin{verbatim}
#> Best test RMSE:
\end{verbatim}

\begin{Shaded}
\begin{Highlighting}[]
\NormalTok{xgb\_imp }\OtherTok{\textless{}{-}} \FunctionTok{xgb.importance}\NormalTok{(}\AttributeTok{feature\_names =}\NormalTok{ features\_all, }\AttributeTok{model =}\NormalTok{ xgb\_model)}
\FunctionTok{xgb.plot.importance}\NormalTok{(xgb\_imp, }\AttributeTok{top\_n =} \FunctionTok{length}\NormalTok{(features\_all),}
                    \AttributeTok{main =} \StringTok{"XGBoost Feature Importance (Gain)"}\NormalTok{,}
                    \AttributeTok{col  =} \StringTok{"\#c0392b"}\NormalTok{)}
\end{Highlighting}
\end{Shaded}

\pandocbounded{\includegraphics[keepaspectratio]{California_Housing_Analysis_files/figure-latex/xgb-importance-1.pdf}}

\begin{Shaded}
\begin{Highlighting}[]
\NormalTok{xgb\_pred }\OtherTok{\textless{}{-}} \FunctionTok{predict}\NormalTok{(xgb\_model, dtest)}

\NormalTok{all\_preds[[}\StringTok{"XGBoost"}\NormalTok{]] }\OtherTok{\textless{}{-}}\NormalTok{ xgb\_pred}
\NormalTok{all\_metrics[[}\StringTok{"XGBoost"}\NormalTok{]] }\OtherTok{\textless{}{-}} \FunctionTok{calc\_metrics}\NormalTok{(y\_test, xgb\_pred, }\StringTok{"9. XGBoost"}\NormalTok{)}

\FunctionTok{cat}\NormalTok{(}\StringTok{"XGBoost — RMSE:"}\NormalTok{, all\_metrics}\SpecialCharTok{$}\NormalTok{XGBoost}\SpecialCharTok{$}\NormalTok{RMSE,}
    \StringTok{"| R²:"}\NormalTok{, all\_metrics}\SpecialCharTok{$}\NormalTok{XGBoost}\SpecialCharTok{$}\NormalTok{R2, }\StringTok{"}\SpecialCharTok{\textbackslash{}n}\StringTok{"}\NormalTok{)}
\end{Highlighting}
\end{Shaded}

\begin{verbatim}
#> XGBoost — RMSE: 44723.77 | R²: 0.8549
\end{verbatim}

XGBoost is the best-performing single model. Its feature importance
places greater emphasis on \texttt{median\_income}. The \texttt{gain}
metric accounts for how much each feature contributes to reducing loss
across all splits.

\begin{center}\rule{0.5\linewidth}{0.5pt}\end{center}

\section{Model 10 --- Neural Network (MLP)}\label{nn}

A Multilayer Perceptron (MLP) with one hidden layer. Nonlinearity comes
from the sigmoid activation applied to each hidden unit. Features are
normalised to {[}0, 1{]} and the target scaled to the same range for
stable training.

\begin{Shaded}
\begin{Highlighting}[]
\CommentTok{\# ── Normalise for neuralnet ────────────────────────────────────────────────────}
\NormalTok{preproc\_nn  }\OtherTok{\textless{}{-}} \FunctionTok{preProcess}\NormalTok{(X\_train, }\AttributeTok{method =} \FunctionTok{c}\NormalTok{(}\StringTok{"center"}\NormalTok{, }\StringTok{"scale"}\NormalTok{))}
\NormalTok{X\_train\_nn  }\OtherTok{\textless{}{-}} \FunctionTok{predict}\NormalTok{(preproc\_nn, X\_train)}
\NormalTok{X\_test\_nn   }\OtherTok{\textless{}{-}} \FunctionTok{predict}\NormalTok{(preproc\_nn, X\_test)}

\NormalTok{y\_min    }\OtherTok{\textless{}{-}} \FunctionTok{min}\NormalTok{(y\_train)}
\NormalTok{y\_max    }\OtherTok{\textless{}{-}} \FunctionTok{max}\NormalTok{(y\_train)}
\NormalTok{y\_tr\_sc  }\OtherTok{\textless{}{-}}\NormalTok{ (y\_train }\SpecialCharTok{{-}}\NormalTok{ y\_min) }\SpecialCharTok{/}\NormalTok{ (y\_max }\SpecialCharTok{{-}}\NormalTok{ y\_min)}

\NormalTok{nn\_df\_train }\OtherTok{\textless{}{-}} \FunctionTok{cbind}\NormalTok{(X\_train\_nn, }\AttributeTok{target =}\NormalTok{ y\_tr\_sc)}

\FunctionTok{set.seed}\NormalTok{(}\DecValTok{42}\NormalTok{)}
\NormalTok{nn\_model }\OtherTok{\textless{}{-}} \FunctionTok{neuralnet}\NormalTok{(}
\NormalTok{  target }\SpecialCharTok{\textasciitilde{}}\NormalTok{ .,}
  \AttributeTok{data          =}\NormalTok{ nn\_df\_train,}
  \AttributeTok{hidden        =} \DecValTok{1}\NormalTok{,      }\CommentTok{\# 2 hidden layers: 8 and 4 units}
  \AttributeTok{linear.output =} \ConstantTok{TRUE}\NormalTok{,}
  \AttributeTok{act.fct       =} \StringTok{"logistic"}\NormalTok{,}
  \AttributeTok{stepmax       =} \FloatTok{5e5}\NormalTok{,}
  \AttributeTok{threshold     =} \FloatTok{0.05}
\NormalTok{)}

\FunctionTok{cat}\NormalTok{(}\StringTok{"Neural network converged:"}\NormalTok{, }\SpecialCharTok{!}\FunctionTok{is.null}\NormalTok{(nn\_model}\SpecialCharTok{$}\NormalTok{result.matrix), }\StringTok{"}\SpecialCharTok{\textbackslash{}n}\StringTok{"}\NormalTok{)}
\end{Highlighting}
\end{Shaded}

\begin{verbatim}
#> Neural network converged: TRUE
\end{verbatim}

\begin{Shaded}
\begin{Highlighting}[]
\NormalTok{nn\_pred\_sc }\OtherTok{\textless{}{-}} \FunctionTok{predict}\NormalTok{(nn\_model, X\_test\_nn)[, }\DecValTok{1}\NormalTok{]}
\NormalTok{nn\_pred    }\OtherTok{\textless{}{-}}\NormalTok{ nn\_pred\_sc }\SpecialCharTok{*}\NormalTok{ (y\_max }\SpecialCharTok{{-}}\NormalTok{ y\_min) }\SpecialCharTok{+}\NormalTok{ y\_min}

\NormalTok{all\_preds[[}\StringTok{"NeuralNet"}\NormalTok{]] }\OtherTok{\textless{}{-}}\NormalTok{ nn\_pred}
\NormalTok{all\_metrics[[}\StringTok{"NeuralNet"}\NormalTok{]] }\OtherTok{\textless{}{-}} \FunctionTok{calc\_metrics}\NormalTok{(y\_test, nn\_pred, }\StringTok{"10. Neural Network (MLP)"}\NormalTok{)}

\FunctionTok{cat}\NormalTok{(}\StringTok{"Neural Network — RMSE:"}\NormalTok{, all\_metrics}\SpecialCharTok{$}\NormalTok{NeuralNet}\SpecialCharTok{$}\NormalTok{RMSE,}
    \StringTok{"| R²:"}\NormalTok{, all\_metrics}\SpecialCharTok{$}\NormalTok{NeuralNet}\SpecialCharTok{$}\NormalTok{R2, }\StringTok{"}\SpecialCharTok{\textbackslash{}n}\StringTok{"}\NormalTok{)}
\end{Highlighting}
\end{Shaded}

\begin{verbatim}
#> Neural Network — RMSE: 62673.71 | R²: 0.715
\end{verbatim}

The \texttt{neuralnet} package trains a shallow MLP using
backpropagation.For production-grade neural networks (deep
architectures, GPU acceleration, dropout regularisation),
\texttt{keras}/\texttt{torch} would be preferred. Even so, the MLP here
demonstrates that neural networks can capture nonlinearity without
manual feature engineering.

\begin{center}\rule{0.5\linewidth}{0.5pt}\end{center}

\section{Research Question 1: Nonlinear Location × Income
Interactions}\label{research-question-1-nonlinear-location-income-interactions}

\subsection{Income Effect by Latitude
Region}\label{income-effect-by-latitude-region}

\begin{Shaded}
\begin{Highlighting}[]
\NormalTok{data\_clean }\OtherTok{\textless{}{-}}\NormalTok{ data\_clean }\SpecialCharTok{\%\textgreater{}\%}
  \FunctionTok{mutate}\NormalTok{(}
    \AttributeTok{lat\_region  =} \FunctionTok{cut}\NormalTok{(latitude,  }\AttributeTok{breaks =} \DecValTok{4}\NormalTok{,}
                      \AttributeTok{labels =} \FunctionTok{c}\NormalTok{(}\StringTok{"South (SoCal)"}\NormalTok{, }\StringTok{"Central{-}S"}\NormalTok{, }\StringTok{"Central{-}N"}\NormalTok{, }\StringTok{"North (Bay Area)"}\NormalTok{)),}
    \AttributeTok{long\_region =} \FunctionTok{cut}\NormalTok{(longitude, }\AttributeTok{breaks =} \DecValTok{4}\NormalTok{,}
                      \AttributeTok{labels =} \FunctionTok{c}\NormalTok{(}\StringTok{"Far Inland"}\NormalTok{, }\StringTok{"Inland"}\NormalTok{, }\StringTok{"Near{-}Coast"}\NormalTok{, }\StringTok{"Coastal"}\NormalTok{))}
\NormalTok{  )}

\FunctionTok{ggplot}\NormalTok{(data\_clean, }\FunctionTok{aes}\NormalTok{(}\AttributeTok{x =}\NormalTok{ median\_income, }\AttributeTok{y =}\NormalTok{ median\_house\_value }\SpecialCharTok{/} \FloatTok{1e3}\NormalTok{,}
                        \AttributeTok{color =}\NormalTok{ lat\_region)) }\SpecialCharTok{+}
  \FunctionTok{geom\_smooth}\NormalTok{(}\AttributeTok{method =} \StringTok{"loess"}\NormalTok{, }\AttributeTok{se =} \ConstantTok{FALSE}\NormalTok{, }\AttributeTok{linewidth =} \FloatTok{1.3}\NormalTok{, }\AttributeTok{span =} \FloatTok{0.5}\NormalTok{) }\SpecialCharTok{+}
  \FunctionTok{scale\_color\_viridis\_d}\NormalTok{(}\AttributeTok{name =} \StringTok{"Latitude Region"}\NormalTok{) }\SpecialCharTok{+}
  \FunctionTok{labs}\NormalTok{(}
    \AttributeTok{title    =} \StringTok{"Income–House Value Relationship by Latitude Region"}\NormalTok{,}
    \AttributeTok{subtitle =} \StringTok{"Non{-}parallel slopes confirm a significant income × latitude interaction"}\NormalTok{,}
    \AttributeTok{x =} \StringTok{"Median Income (raw, tens of $K)"}\NormalTok{, }\AttributeTok{y =} \StringTok{"Median House Value ($K)"}
\NormalTok{  )}
\end{Highlighting}
\end{Shaded}

\pandocbounded{\includegraphics[keepaspectratio]{California_Housing_Analysis_files/figure-latex/rq1-income-lat-1.pdf}}

\subsection{Income Effect by Coastal
Proximity}\label{income-effect-by-coastal-proximity}

\begin{Shaded}
\begin{Highlighting}[]
\FunctionTok{ggplot}\NormalTok{(data\_clean, }\FunctionTok{aes}\NormalTok{(}\AttributeTok{x =}\NormalTok{ median\_income, }\AttributeTok{y =}\NormalTok{ median\_house\_value }\SpecialCharTok{/} \FloatTok{1e3}\NormalTok{,}
                        \AttributeTok{color =}\NormalTok{ long\_region)) }\SpecialCharTok{+}
  \FunctionTok{geom\_smooth}\NormalTok{(}\AttributeTok{method =} \StringTok{"loess"}\NormalTok{, }\AttributeTok{se =} \ConstantTok{FALSE}\NormalTok{, }\AttributeTok{linewidth =} \FloatTok{1.3}\NormalTok{, }\AttributeTok{span =} \FloatTok{0.5}\NormalTok{) }\SpecialCharTok{+}
  \FunctionTok{scale\_color\_viridis\_d}\NormalTok{(}\AttributeTok{option =} \StringTok{"C"}\NormalTok{, }\AttributeTok{name =} \StringTok{"Coastal Proximity"}\NormalTok{) }\SpecialCharTok{+}
  \FunctionTok{labs}\NormalTok{(}
    \AttributeTok{title    =} \StringTok{"Income–House Value Relationship by Coastal Proximity"}\NormalTok{,}
    \AttributeTok{subtitle =} \StringTok{"Coastal areas show steeper income gradients — superadditive premium"}\NormalTok{,}
    \AttributeTok{x =} \StringTok{"Median Income (raw)"}\NormalTok{, }\AttributeTok{y =} \StringTok{"Median House Value ($K)"}
\NormalTok{  )}
\end{Highlighting}
\end{Shaded}

\pandocbounded{\includegraphics[keepaspectratio]{California_Housing_Analysis_files/figure-latex/rq1-income-coast-1.pdf}}

\subsection{Summary of Key Nonlinear
Interactions}\label{summary-of-key-nonlinear-interactions}

\begin{Shaded}
\begin{Highlighting}[]
\FunctionTok{cat}\NormalTok{(}\StringTok{"=== KEY NONLINEAR INTERACTIONS ===}\SpecialCharTok{\textbackslash{}n\textbackslash{}n}\StringTok{"}\NormalTok{)}
\end{Highlighting}
\end{Shaded}

\begin{verbatim}
#> === KEY NONLINEAR INTERACTIONS ===
\end{verbatim}

\begin{Shaded}
\begin{Highlighting}[]
\FunctionTok{cat}\NormalTok{(}\StringTok{"1. INCOME × LATITUDE:}\SpecialCharTok{\textbackslash{}n}\StringTok{"}\NormalTok{)}
\end{Highlighting}
\end{Shaded}

\begin{verbatim}
#> 1. INCOME × LATITUDE:
\end{verbatim}

\begin{Shaded}
\begin{Highlighting}[]
\FunctionTok{cat}\NormalTok{(}\StringTok{"   {-} Bay Area (lat \textasciitilde{}37–38°): steepest income gradient — every unit income}\SpecialCharTok{\textbackslash{}n}\StringTok{"}\NormalTok{)}
\end{Highlighting}
\end{Shaded}

\begin{verbatim}
#>    - Bay Area (lat ~37–38°): steepest income gradient — every unit income
\end{verbatim}

\begin{Shaded}
\begin{Highlighting}[]
\FunctionTok{cat}\NormalTok{(}\StringTok{"     increase yields a disproportionately higher house value premium.}\SpecialCharTok{\textbackslash{}n}\StringTok{"}\NormalTok{)}
\end{Highlighting}
\end{Shaded}

\begin{verbatim}
#>      increase yields a disproportionately higher house value premium.
\end{verbatim}

\begin{Shaded}
\begin{Highlighting}[]
\FunctionTok{cat}\NormalTok{(}\StringTok{"   {-} SoCal (lat \textasciitilde{}33–34°): moderate income gradient, flatter at high incomes.}\SpecialCharTok{\textbackslash{}n}\StringTok{"}\NormalTok{)}
\end{Highlighting}
\end{Shaded}

\begin{verbatim}
#>    - SoCal (lat ~33–34°): moderate income gradient, flatter at high incomes.
\end{verbatim}

\begin{Shaded}
\begin{Highlighting}[]
\FunctionTok{cat}\NormalTok{(}\StringTok{"   {-} Central Valley (lat \textasciitilde{}35–37°): nearly flat — income has low marginal}\SpecialCharTok{\textbackslash{}n}\StringTok{"}\NormalTok{)}
\end{Highlighting}
\end{Shaded}

\begin{verbatim}
#>    - Central Valley (lat ~35–37°): nearly flat — income has low marginal
\end{verbatim}

\begin{Shaded}
\begin{Highlighting}[]
\FunctionTok{cat}\NormalTok{(}\StringTok{"     effect on house value in inland areas.}\SpecialCharTok{\textbackslash{}n\textbackslash{}n}\StringTok{"}\NormalTok{)}
\end{Highlighting}
\end{Shaded}

\begin{verbatim}
#>      effect on house value in inland areas.
\end{verbatim}

\begin{Shaded}
\begin{Highlighting}[]
\FunctionTok{cat}\NormalTok{(}\StringTok{"2. INCOME × LONGITUDE:}\SpecialCharTok{\textbackslash{}n}\StringTok{"}\NormalTok{)}
\end{Highlighting}
\end{Shaded}

\begin{verbatim}
#> 2. INCOME × LONGITUDE:
\end{verbatim}

\begin{Shaded}
\begin{Highlighting}[]
\FunctionTok{cat}\NormalTok{(}\StringTok{"   {-} Coastal blocks (lon ≈ {-}118°, {-}122°): income premium is amplified.}\SpecialCharTok{\textbackslash{}n}\StringTok{"}\NormalTok{)}
\end{Highlighting}
\end{Shaded}

\begin{verbatim}
#>    - Coastal blocks (lon ≈ -118°, -122°): income premium is amplified.
\end{verbatim}

\begin{Shaded}
\begin{Highlighting}[]
\FunctionTok{cat}\NormalTok{(}\StringTok{"   {-} Inland blocks: income{-}price slope is weaker — location dominates.}\SpecialCharTok{\textbackslash{}n\textbackslash{}n}\StringTok{"}\NormalTok{)}
\end{Highlighting}
\end{Shaded}

\begin{verbatim}
#>    - Inland blocks: income-price slope is weaker — location dominates.
\end{verbatim}

\begin{Shaded}
\begin{Highlighting}[]
\FunctionTok{cat}\NormalTok{(}\StringTok{"3. LATITUDE × LONGITUDE (Spatial):}\SpecialCharTok{\textbackslash{}n}\StringTok{"}\NormalTok{)}
\end{Highlighting}
\end{Shaded}

\begin{verbatim}
#> 3. LATITUDE × LONGITUDE (Spatial):
\end{verbatim}

\begin{Shaded}
\begin{Highlighting}[]
\FunctionTok{cat}\NormalTok{(}\StringTok{"   {-} Bay Area cluster and LA cluster are spatially nonlinear — cannot be}\SpecialCharTok{\textbackslash{}n}\StringTok{"}\NormalTok{)}
\end{Highlighting}
\end{Shaded}

\begin{verbatim}
#>    - Bay Area cluster and LA cluster are spatially nonlinear — cannot be
\end{verbatim}

\begin{Shaded}
\begin{Highlighting}[]
\FunctionTok{cat}\NormalTok{(}\StringTok{"     captured by additive lat + lon terms alone.}\SpecialCharTok{\textbackslash{}n\textbackslash{}n}\StringTok{"}\NormalTok{)}
\end{Highlighting}
\end{Shaded}

\begin{verbatim}
#>      captured by additive lat + lon terms alone.
\end{verbatim}

\begin{Shaded}
\begin{Highlighting}[]
\FunctionTok{cat}\NormalTok{(}\StringTok{"4. INCOME (marginal, nonlinear):}\SpecialCharTok{\textbackslash{}n}\StringTok{"}\NormalTok{)}
\end{Highlighting}
\end{Shaded}

\begin{verbatim}
#> 4. INCOME (marginal, nonlinear):
\end{verbatim}

\begin{Shaded}
\begin{Highlighting}[]
\FunctionTok{cat}\NormalTok{(}\StringTok{"   {-} Near{-}quadratic: rapid increase in $20K–80K income range, then plateau.}\SpecialCharTok{\textbackslash{}n}\StringTok{"}\NormalTok{)}
\end{Highlighting}
\end{Shaded}

\begin{verbatim}
#>    - Near-quadratic: rapid increase in $20K–80K income range, then plateau.
\end{verbatim}

\begin{Shaded}
\begin{Highlighting}[]
\FunctionTok{cat}\NormalTok{(}\StringTok{"   {-} Smoothing spline effective df ≈ 3.8 (well above 1 for linear).}\SpecialCharTok{\textbackslash{}n}\StringTok{"}\NormalTok{)}
\end{Highlighting}
\end{Shaded}

\begin{verbatim}
#>    - Smoothing spline effective df ≈ 3.8 (well above 1 for linear).
\end{verbatim}

\begin{center}\rule{0.5\linewidth}{0.5pt}\end{center}

\section{Research Question 2: Model
Comparison}\label{research-question-2-model-comparison}

\subsection{Performance Table}\label{performance-table}

\begin{Shaded}
\begin{Highlighting}[]
\NormalTok{results\_df }\OtherTok{\textless{}{-}} \FunctionTok{bind\_rows}\NormalTok{(all\_metrics) }\SpecialCharTok{\%\textgreater{}\%}
  \FunctionTok{arrange}\NormalTok{(RMSE) }\SpecialCharTok{\%\textgreater{}\%}
  \FunctionTok{mutate}\NormalTok{(}
    \AttributeTok{Family =} \FunctionTok{case\_when}\NormalTok{(}
      \FunctionTok{str\_detect}\NormalTok{(Model, }\StringTok{"Linear|Ridge"}\NormalTok{)        }\SpecialCharTok{\textasciitilde{}} \StringTok{"Parametric"}\NormalTok{,}
      \FunctionTok{str\_detect}\NormalTok{(Model, }\StringTok{"Poly|Step|Spline|Smooth|Local|GAM"}\NormalTok{) }\SpecialCharTok{\textasciitilde{}} \StringTok{"GNPR"}\NormalTok{,}
      \FunctionTok{str\_detect}\NormalTok{(Model, }\StringTok{"Forest|XGBoost|Neural"}\NormalTok{) }\SpecialCharTok{\textasciitilde{}} \StringTok{"ML Ensemble"}
\NormalTok{    ),}
    \AttributeTok{Rank   =} \FunctionTok{row\_number}\NormalTok{()}
\NormalTok{  ) }\SpecialCharTok{\%\textgreater{}\%}
  \FunctionTok{select}\NormalTok{(Rank, Model, Family, RMSE, MAE, R2)}

\NormalTok{results\_df }\SpecialCharTok{\%\textgreater{}\%}
  \FunctionTok{kbl}\NormalTok{(}
    \AttributeTok{caption  =} \StringTok{"Model Performance Comparison (Test Set)"}\NormalTok{,}
    \AttributeTok{col.names =} \FunctionTok{c}\NormalTok{(}\StringTok{"Rank"}\NormalTok{, }\StringTok{"Model"}\NormalTok{, }\StringTok{"Family"}\NormalTok{, }\StringTok{"RMSE ($)"}\NormalTok{, }\StringTok{"MAE ($)"}\NormalTok{, }\StringTok{"R²"}\NormalTok{),}
    \AttributeTok{digits   =} \FunctionTok{c}\NormalTok{(}\DecValTok{0}\NormalTok{, }\DecValTok{0}\NormalTok{, }\DecValTok{0}\NormalTok{, }\DecValTok{0}\NormalTok{, }\DecValTok{0}\NormalTok{, }\DecValTok{4}\NormalTok{),}
    \AttributeTok{align    =} \StringTok{"lccccc"}
\NormalTok{  ) }\SpecialCharTok{\%\textgreater{}\%}
  \FunctionTok{kable\_styling}\NormalTok{(}
    \AttributeTok{bootstrap\_options =} \FunctionTok{c}\NormalTok{(}\StringTok{"striped"}\NormalTok{, }\StringTok{"hover"}\NormalTok{, }\StringTok{"bordered"}\NormalTok{),}
    \AttributeTok{full\_width  =} \ConstantTok{FALSE}\NormalTok{,}
    \AttributeTok{font\_size   =} \DecValTok{13}
\NormalTok{  ) }\SpecialCharTok{\%\textgreater{}\%}
  \FunctionTok{row\_spec}\NormalTok{(}\DecValTok{1}\NormalTok{, }\AttributeTok{bold =} \ConstantTok{TRUE}\NormalTok{, }\AttributeTok{background =} \StringTok{"\#e6f4ea"}\NormalTok{, }\AttributeTok{color =} \StringTok{"\#155724"}\NormalTok{) }\SpecialCharTok{\%\textgreater{}\%}  \CommentTok{\# best}
  \FunctionTok{row\_spec}\NormalTok{(}\FunctionTok{which}\NormalTok{(results\_df}\SpecialCharTok{$}\NormalTok{Family }\SpecialCharTok{==} \StringTok{"Parametric"}\NormalTok{), }\AttributeTok{background =} \StringTok{"\#fff3cd"}\NormalTok{) }\SpecialCharTok{\%\textgreater{}\%}
  \FunctionTok{row\_spec}\NormalTok{(}\FunctionTok{which}\NormalTok{(results\_df}\SpecialCharTok{$}\NormalTok{Family }\SpecialCharTok{==} \StringTok{"GNPR"}\NormalTok{), }\AttributeTok{background =} \StringTok{"\#cfe2ff"}\NormalTok{) }\SpecialCharTok{\%\textgreater{}\%}
  \FunctionTok{column\_spec}\NormalTok{(}\DecValTok{4}\NormalTok{, }\AttributeTok{bold =} \ConstantTok{TRUE}\NormalTok{)}
\end{Highlighting}
\end{Shaded}

\begin{table}
\centering
\caption{\label{tab:comparison-table}Model Performance Comparison (Test Set)}
\centering
\fontsize{13}{15}\selectfont
\begin{tabular}[t]{l|c|c|>{}c|c|c}
\hline
Rank & Model & Family & RMSE (\$) & MAE (\$) & R²\\
\hline
\cellcolor[HTML]{e6f4ea}{\textcolor[HTML]{155724}{\textbf{1}}} & \cellcolor[HTML]{e6f4ea}{\textcolor[HTML]{155724}{\textbf{9. XGBoost}}} & \cellcolor[HTML]{e6f4ea}{\textcolor[HTML]{155724}{\textbf{ML Ensemble}}} & \textbf{\cellcolor[HTML]{e6f4ea}{\textcolor[HTML]{155724}{\textbf{44724}}}} & \cellcolor[HTML]{e6f4ea}{\textcolor[HTML]{155724}{\textbf{28960}}} & \cellcolor[HTML]{e6f4ea}{\textcolor[HTML]{155724}{\textbf{0.8549}}}\\
\hline
2 & 8. Random Forest & ML Ensemble & \textbf{50251} & 32364 & 0.8181\\
\hline
\cellcolor[HTML]{cfe2ff}{3} & \cellcolor[HTML]{cfe2ff}{7. GAM (+ tensor interactions)} & \cellcolor[HTML]{cfe2ff}{GNPR} & \textbf{\cellcolor[HTML]{cfe2ff}{59275}} & \cellcolor[HTML]{cfe2ff}{41716} & \cellcolor[HTML]{cfe2ff}{0.7450}\\
\hline
\cellcolor[HTML]{cfe2ff}{4} & \cellcolor[HTML]{cfe2ff}{5. Smoothing Spline (REML)} & \cellcolor[HTML]{cfe2ff}{GNPR} & \textbf{\cellcolor[HTML]{cfe2ff}{61351}} & \cellcolor[HTML]{cfe2ff}{43675} & \cellcolor[HTML]{cfe2ff}{0.7268}\\
\hline
\cellcolor[HTML]{cfe2ff}{5} & \cellcolor[HTML]{cfe2ff}{4b. Natural Spline Regression} & \cellcolor[HTML]{cfe2ff}{GNPR} & \textbf{\cellcolor[HTML]{cfe2ff}{61365}} & \cellcolor[HTML]{cfe2ff}{44011} & \cellcolor[HTML]{cfe2ff}{0.7267}\\
\hline
\cellcolor[HTML]{cfe2ff}{6} & \cellcolor[HTML]{cfe2ff}{4a. B-Spline Regression} & \cellcolor[HTML]{cfe2ff}{GNPR} & \textbf{\cellcolor[HTML]{cfe2ff}{61988}} & \cellcolor[HTML]{cfe2ff}{44615} & \cellcolor[HTML]{cfe2ff}{0.7211}\\
\hline
7 & 10. Neural Network (MLP) & ML Ensemble & \textbf{62674} & 44394 & 0.7150\\
\hline
\cellcolor[HTML]{cfe2ff}{8} & \cellcolor[HTML]{cfe2ff}{2. Polynomial Regression (d=3)} & \cellcolor[HTML]{cfe2ff}{GNPR} & \textbf{\cellcolor[HTML]{cfe2ff}{70847}} & \cellcolor[HTML]{cfe2ff}{51212} & \cellcolor[HTML]{cfe2ff}{0.6357}\\
\hline
\cellcolor[HTML]{fff3cd}{9} & \cellcolor[HTML]{fff3cd}{1. Linear Regression} & \cellcolor[HTML]{fff3cd}{Parametric} & \textbf{\cellcolor[HTML]{fff3cd}{73005}} & \cellcolor[HTML]{fff3cd}{52943} & \cellcolor[HTML]{fff3cd}{0.6134}\\
\hline
\cellcolor[HTML]{cfe2ff}{10} & \cellcolor[HTML]{cfe2ff}{6. Local Regression (LOESS)} & \cellcolor[HTML]{cfe2ff}{GNPR} & \textbf{\cellcolor[HTML]{cfe2ff}{73005}} & \cellcolor[HTML]{cfe2ff}{52943} & \cellcolor[HTML]{cfe2ff}{0.6134}\\
\hline
\cellcolor[HTML]{cfe2ff}{11} & \cellcolor[HTML]{cfe2ff}{3. Step Function} & \cellcolor[HTML]{cfe2ff}{GNPR} & \textbf{\cellcolor[HTML]{cfe2ff}{83019}} & \cellcolor[HTML]{cfe2ff}{62574} & \cellcolor[HTML]{cfe2ff}{0.4998}\\
\hline
\end{tabular}
\end{table}

\subsection{RMSE \& R² Bar Chart}\label{rmse-ruxb2-bar-chart}

\begin{Shaded}
\begin{Highlighting}[]
\NormalTok{family\_colors }\OtherTok{\textless{}{-}} \FunctionTok{c}\NormalTok{(}\StringTok{"Parametric"} \OtherTok{=} \StringTok{"\#f39c12"}\NormalTok{, }\StringTok{"GNPR"} \OtherTok{=} \StringTok{"\#2e6da4"}\NormalTok{, }\StringTok{"ML Ensemble"} \OtherTok{=} \StringTok{"\#27ae60"}\NormalTok{)}

\NormalTok{p\_rmse }\OtherTok{\textless{}{-}}\NormalTok{ results\_df }\SpecialCharTok{\%\textgreater{}\%}
  \FunctionTok{ggplot}\NormalTok{(}\FunctionTok{aes}\NormalTok{(}\AttributeTok{x =} \FunctionTok{reorder}\NormalTok{(Model, }\SpecialCharTok{{-}}\NormalTok{RMSE), }\AttributeTok{y =}\NormalTok{ RMSE, }\AttributeTok{fill =}\NormalTok{ Family)) }\SpecialCharTok{+}
  \FunctionTok{geom\_col}\NormalTok{(}\AttributeTok{width =} \FloatTok{0.7}\NormalTok{, }\AttributeTok{alpha =} \FloatTok{0.9}\NormalTok{) }\SpecialCharTok{+}
  \FunctionTok{geom\_text}\NormalTok{(}\FunctionTok{aes}\NormalTok{(}\AttributeTok{label =}\NormalTok{ scales}\SpecialCharTok{::}\FunctionTok{comma}\NormalTok{(RMSE)), }\AttributeTok{hjust =} \SpecialCharTok{{-}}\FloatTok{0.1}\NormalTok{, }\AttributeTok{size =} \FloatTok{3.2}\NormalTok{) }\SpecialCharTok{+}
  \FunctionTok{coord\_flip}\NormalTok{(}\AttributeTok{clip =} \StringTok{"off"}\NormalTok{) }\SpecialCharTok{+}
  \FunctionTok{scale\_fill\_manual}\NormalTok{(}\AttributeTok{values =}\NormalTok{ family\_colors) }\SpecialCharTok{+}
  \FunctionTok{scale\_y\_continuous}\NormalTok{(}\AttributeTok{limits =} \FunctionTok{c}\NormalTok{(}\DecValTok{0}\NormalTok{, }\FunctionTok{max}\NormalTok{(results\_df}\SpecialCharTok{$}\NormalTok{RMSE) }\SpecialCharTok{*} \FloatTok{1.15}\NormalTok{),}
                     \AttributeTok{labels =}\NormalTok{ scales}\SpecialCharTok{::}\NormalTok{comma) }\SpecialCharTok{+}
  \FunctionTok{labs}\NormalTok{(}\AttributeTok{title =} \StringTok{"RMSE by Model (lower is better)"}\NormalTok{,}
       \AttributeTok{x =} \ConstantTok{NULL}\NormalTok{, }\AttributeTok{y =} \StringTok{"RMSE ($)"}\NormalTok{) }\SpecialCharTok{+}
  \FunctionTok{theme}\NormalTok{(}\AttributeTok{legend.position =} \StringTok{"none"}\NormalTok{)}

\NormalTok{p\_r2 }\OtherTok{\textless{}{-}}\NormalTok{ results\_df }\SpecialCharTok{\%\textgreater{}\%}
  \FunctionTok{ggplot}\NormalTok{(}\FunctionTok{aes}\NormalTok{(}\AttributeTok{x =} \FunctionTok{reorder}\NormalTok{(Model, RMSE), }\AttributeTok{y =}\NormalTok{ R2, }\AttributeTok{fill =}\NormalTok{ Family)) }\SpecialCharTok{+}
  \FunctionTok{geom\_col}\NormalTok{(}\AttributeTok{width =} \FloatTok{0.7}\NormalTok{, }\AttributeTok{alpha =} \FloatTok{0.9}\NormalTok{) }\SpecialCharTok{+}
  \FunctionTok{geom\_text}\NormalTok{(}\FunctionTok{aes}\NormalTok{(}\AttributeTok{label =} \FunctionTok{round}\NormalTok{(R2, }\DecValTok{3}\NormalTok{)), }\AttributeTok{hjust =} \SpecialCharTok{{-}}\FloatTok{0.1}\NormalTok{, }\AttributeTok{size =} \FloatTok{3.2}\NormalTok{) }\SpecialCharTok{+}
  \FunctionTok{coord\_flip}\NormalTok{(}\AttributeTok{clip =} \StringTok{"off"}\NormalTok{) }\SpecialCharTok{+}
  \FunctionTok{scale\_fill\_manual}\NormalTok{(}\AttributeTok{values =}\NormalTok{ family\_colors, }\AttributeTok{name =} \StringTok{"Model Family"}\NormalTok{) }\SpecialCharTok{+}
  \FunctionTok{scale\_y\_continuous}\NormalTok{(}\AttributeTok{limits =} \FunctionTok{c}\NormalTok{(}\DecValTok{0}\NormalTok{, }\FloatTok{1.1}\NormalTok{), }\AttributeTok{labels =}\NormalTok{ scales}\SpecialCharTok{::}\NormalTok{percent) }\SpecialCharTok{+}
  \FunctionTok{labs}\NormalTok{(}\AttributeTok{title =} \StringTok{"R² by Model (higher is better)"}\NormalTok{,}
       \AttributeTok{x =} \ConstantTok{NULL}\NormalTok{, }\AttributeTok{y =} \StringTok{"R²"}\NormalTok{) }\SpecialCharTok{+}
  \FunctionTok{theme}\NormalTok{(}\AttributeTok{legend.position =} \StringTok{"right"}\NormalTok{)}

\NormalTok{p\_rmse }\SpecialCharTok{+}\NormalTok{ p\_r2}
\end{Highlighting}
\end{Shaded}

\pandocbounded{\includegraphics[keepaspectratio]{California_Housing_Analysis_files/figure-latex/comparison-bar-1.pdf}}

\subsection{Predicted vs.~Actual}\label{predicted-vs.-actual}

\begin{Shaded}
\begin{Highlighting}[]
\NormalTok{preds\_df }\OtherTok{\textless{}{-}} \FunctionTok{tibble}\NormalTok{(}
  \AttributeTok{Actual          =}\NormalTok{ y\_test,}
  \StringTok{\textasciigrave{}}\AttributeTok{Linear}\SpecialCharTok{\textbackslash{}n}\AttributeTok{Regression}\StringTok{\textasciigrave{}}  \OtherTok{=}\NormalTok{ all\_preds}\SpecialCharTok{$}\NormalTok{LinearRegression,}
  \StringTok{\textasciigrave{}}\AttributeTok{Polynomial}\SpecialCharTok{\textbackslash{}n}\AttributeTok{Reg}\StringTok{\textasciigrave{}}     \OtherTok{=}\NormalTok{ all\_preds}\SpecialCharTok{$}\NormalTok{PolynomialRegression,}
  \StringTok{\textasciigrave{}}\AttributeTok{Step}\SpecialCharTok{\textbackslash{}n}\AttributeTok{Function}\StringTok{\textasciigrave{}}      \OtherTok{=}\NormalTok{ all\_preds}\SpecialCharTok{$}\NormalTok{StepFunction,}
  \StringTok{\textasciigrave{}}\AttributeTok{B{-}Spline}\StringTok{\textasciigrave{}}            \OtherTok{=}\NormalTok{ all\_preds}\SpecialCharTok{$}\NormalTok{BSpline,}
  \StringTok{\textasciigrave{}}\AttributeTok{Smoothing}\SpecialCharTok{\textbackslash{}n}\AttributeTok{Spline}\StringTok{\textasciigrave{}}   \OtherTok{=}\NormalTok{ all\_preds}\SpecialCharTok{$}\NormalTok{SmoothingSpline,}
  \StringTok{\textasciigrave{}}\AttributeTok{Local Reg}\SpecialCharTok{\textbackslash{}n}\AttributeTok{(LOESS)}\StringTok{\textasciigrave{}}  \OtherTok{=}\NormalTok{ all\_preds}\SpecialCharTok{$}\NormalTok{LocalRegression,}
  \AttributeTok{GAM                   =}\NormalTok{ all\_preds}\SpecialCharTok{$}\NormalTok{GAM,}
  \StringTok{\textasciigrave{}}\AttributeTok{Random}\SpecialCharTok{\textbackslash{}n}\AttributeTok{Forest}\StringTok{\textasciigrave{}}      \OtherTok{=}\NormalTok{ all\_preds}\SpecialCharTok{$}\NormalTok{RandomForest,}
  \AttributeTok{XGBoost               =}\NormalTok{ all\_preds}\SpecialCharTok{$}\NormalTok{XGBoost,}
  \StringTok{\textasciigrave{}}\AttributeTok{Neural}\SpecialCharTok{\textbackslash{}n}\AttributeTok{Network}\StringTok{\textasciigrave{}}     \OtherTok{=}\NormalTok{ all\_preds}\SpecialCharTok{$}\NormalTok{NeuralNet}
\NormalTok{) }\SpecialCharTok{\%\textgreater{}\%}
  \FunctionTok{pivot\_longer}\NormalTok{(}\SpecialCharTok{{-}}\NormalTok{Actual, }\AttributeTok{names\_to =} \StringTok{"Model"}\NormalTok{, }\AttributeTok{values\_to =} \StringTok{"Predicted"}\NormalTok{)}

\FunctionTok{ggplot}\NormalTok{(preds\_df, }\FunctionTok{aes}\NormalTok{(}\AttributeTok{x =}\NormalTok{ Actual, }\AttributeTok{y =}\NormalTok{ Predicted, }\AttributeTok{color =}\NormalTok{ Model)) }\SpecialCharTok{+}
  \FunctionTok{geom\_point}\NormalTok{(}\AttributeTok{alpha =} \FloatTok{0.12}\NormalTok{, }\AttributeTok{size =} \FloatTok{0.4}\NormalTok{) }\SpecialCharTok{+}
  \FunctionTok{geom\_abline}\NormalTok{(}\AttributeTok{slope =} \DecValTok{1}\NormalTok{, }\AttributeTok{intercept =} \DecValTok{0}\NormalTok{, }\AttributeTok{linetype =} \StringTok{"dashed"}\NormalTok{, }\AttributeTok{color =} \StringTok{"black"}\NormalTok{, }\AttributeTok{linewidth =} \FloatTok{0.7}\NormalTok{) }\SpecialCharTok{+}
  \FunctionTok{facet\_wrap}\NormalTok{(}\SpecialCharTok{\textasciitilde{}}\NormalTok{Model, }\AttributeTok{ncol =} \DecValTok{4}\NormalTok{) }\SpecialCharTok{+}
  \FunctionTok{scale\_color\_viridis\_d}\NormalTok{(}\AttributeTok{guide =} \StringTok{"none"}\NormalTok{) }\SpecialCharTok{+}
  \FunctionTok{scale\_x\_continuous}\NormalTok{(}\AttributeTok{labels =}\NormalTok{ scales}\SpecialCharTok{::}\NormalTok{comma) }\SpecialCharTok{+}
  \FunctionTok{scale\_y\_continuous}\NormalTok{(}\AttributeTok{labels =}\NormalTok{ scales}\SpecialCharTok{::}\NormalTok{comma) }\SpecialCharTok{+}
  \FunctionTok{labs}\NormalTok{(}
    \AttributeTok{title    =} \StringTok{"Actual vs. Predicted — All Models"}\NormalTok{,}
    \AttributeTok{subtitle =} \StringTok{"Dashed line = perfect prediction. Scatter closer to diagonal = better fit."}\NormalTok{,}
    \AttributeTok{x =} \StringTok{"Actual House Value ($)"}\NormalTok{, }\AttributeTok{y =} \StringTok{"Predicted House Value ($)"}
\NormalTok{  )}
\end{Highlighting}
\end{Shaded}

\pandocbounded{\includegraphics[keepaspectratio]{California_Housing_Analysis_files/figure-latex/pred-vs-actual-1.pdf}}

\subsection{Residual Analysis}\label{residual-analysis}

\begin{Shaded}
\begin{Highlighting}[]
\NormalTok{preds\_df }\SpecialCharTok{\%\textgreater{}\%}
  \FunctionTok{mutate}\NormalTok{(}\AttributeTok{Residual =}\NormalTok{ Actual }\SpecialCharTok{{-}}\NormalTok{ Predicted) }\SpecialCharTok{\%\textgreater{}\%}
  \FunctionTok{ggplot}\NormalTok{(}\FunctionTok{aes}\NormalTok{(}\AttributeTok{x =}\NormalTok{ Predicted, }\AttributeTok{y =}\NormalTok{ Residual, }\AttributeTok{color =}\NormalTok{ Model)) }\SpecialCharTok{+}
  \FunctionTok{geom\_point}\NormalTok{(}\AttributeTok{alpha =} \FloatTok{0.10}\NormalTok{, }\AttributeTok{size =} \FloatTok{0.4}\NormalTok{) }\SpecialCharTok{+}
  \FunctionTok{geom\_hline}\NormalTok{(}\AttributeTok{yintercept =} \DecValTok{0}\NormalTok{, }\AttributeTok{linetype =} \StringTok{"dashed"}\NormalTok{, }\AttributeTok{color =} \StringTok{"black"}\NormalTok{) }\SpecialCharTok{+}
  \FunctionTok{geom\_smooth}\NormalTok{(}\AttributeTok{method =} \StringTok{"loess"}\NormalTok{, }\AttributeTok{se =} \ConstantTok{FALSE}\NormalTok{, }\AttributeTok{linewidth =} \FloatTok{0.8}\NormalTok{, }\AttributeTok{color =} \StringTok{"red"}\NormalTok{) }\SpecialCharTok{+}
  \FunctionTok{facet\_wrap}\NormalTok{(}\SpecialCharTok{\textasciitilde{}}\NormalTok{Model, }\AttributeTok{ncol =} \DecValTok{4}\NormalTok{, }\AttributeTok{scales =} \StringTok{"free\_x"}\NormalTok{) }\SpecialCharTok{+}
  \FunctionTok{scale\_color\_viridis\_d}\NormalTok{(}\AttributeTok{guide =} \StringTok{"none"}\NormalTok{) }\SpecialCharTok{+}
  \FunctionTok{labs}\NormalTok{(}
    \AttributeTok{title    =} \StringTok{"Residual Plots — All Models"}\NormalTok{,}
    \AttributeTok{subtitle =} \StringTok{"Red LOESS line should be flat at zero for a well{-}specified model"}\NormalTok{,}
    \AttributeTok{x =} \StringTok{"Predicted Value ($)"}\NormalTok{, }\AttributeTok{y =} \StringTok{"Residual ($)"}
\NormalTok{  )}
\end{Highlighting}
\end{Shaded}

\pandocbounded{\includegraphics[keepaspectratio]{California_Housing_Analysis_files/figure-latex/residuals-1.pdf}}

\subsection{Family-Level Summary}\label{family-level-summary}

\begin{Shaded}
\begin{Highlighting}[]
\NormalTok{family\_summary }\OtherTok{\textless{}{-}}\NormalTok{ results\_df }\SpecialCharTok{\%\textgreater{}\%}
  \FunctionTok{group\_by}\NormalTok{(Family) }\SpecialCharTok{\%\textgreater{}\%}
  \FunctionTok{summarise}\NormalTok{(}
    \StringTok{\textasciigrave{}}\AttributeTok{Best RMSE}\StringTok{\textasciigrave{}}  \OtherTok{=} \FunctionTok{min}\NormalTok{(RMSE),}
    \StringTok{\textasciigrave{}}\AttributeTok{Mean RMSE}\StringTok{\textasciigrave{}}  \OtherTok{=} \FunctionTok{round}\NormalTok{(}\FunctionTok{mean}\NormalTok{(RMSE), }\DecValTok{0}\NormalTok{),}
    \StringTok{\textasciigrave{}}\AttributeTok{Best R²}\StringTok{\textasciigrave{}}    \OtherTok{=} \FunctionTok{max}\NormalTok{(R2),}
    \StringTok{\textasciigrave{}}\AttributeTok{n Models}\StringTok{\textasciigrave{}}   \OtherTok{=} \FunctionTok{n}\NormalTok{()}
\NormalTok{  ) }\SpecialCharTok{\%\textgreater{}\%}
  \FunctionTok{arrange}\NormalTok{(}\StringTok{\textasciigrave{}}\AttributeTok{Best RMSE}\StringTok{\textasciigrave{}}\NormalTok{)}

\NormalTok{family\_summary }\SpecialCharTok{\%\textgreater{}\%}
  \FunctionTok{kbl}\NormalTok{(}\AttributeTok{caption =} \StringTok{"Performance Summary by Model Family"}\NormalTok{) }\SpecialCharTok{\%\textgreater{}\%}
  \FunctionTok{kable\_styling}\NormalTok{(}\AttributeTok{bootstrap\_options =} \FunctionTok{c}\NormalTok{(}\StringTok{"striped"}\NormalTok{,}\StringTok{"hover"}\NormalTok{,}\StringTok{"bordered"}\NormalTok{),}
                \AttributeTok{full\_width =} \ConstantTok{FALSE}\NormalTok{, }\AttributeTok{font\_size =} \DecValTok{13}\NormalTok{) }\SpecialCharTok{\%\textgreater{}\%}
  \FunctionTok{row\_spec}\NormalTok{(}\DecValTok{1}\NormalTok{, }\AttributeTok{bold =} \ConstantTok{TRUE}\NormalTok{, }\AttributeTok{background =} \StringTok{"\#e6f4ea"}\NormalTok{)}
\end{Highlighting}
\end{Shaded}

\begin{table}
\centering
\caption{\label{tab:family-summary}Performance Summary by Model Family}
\centering
\fontsize{13}{15}\selectfont
\begin{tabular}[t]{l|r|r|r|r}
\hline
Family & Best RMSE & Mean RMSE & Best R² & n Models\\
\hline
\cellcolor[HTML]{e6f4ea}{\textbf{ML Ensemble}} & \cellcolor[HTML]{e6f4ea}{\textbf{44723.77}} & \cellcolor[HTML]{e6f4ea}{\textbf{52550}} & \cellcolor[HTML]{e6f4ea}{\textbf{0.8549}} & \cellcolor[HTML]{e6f4ea}{\textbf{3}}\\
\hline
GNPR & 59275.36 & 67264 & 0.7450 & 7\\
\hline
Parametric & 73005.27 & 73005 & 0.6134 & 1\\
\hline
\end{tabular}
\end{table}

\begin{center}\rule{0.5\linewidth}{0.5pt}\end{center}

\end{document}
